\documentclass[11pt,a4paper]{article}
\usepackage[T1]{fontenc}
\usepackage[utf8]{inputenc}
\usepackage{lmodern,amsmath,amssymb,amsthm,mathtools}
\usepackage[a4paper,margin=26mm,headheight=15pt]{geometry}
\usepackage{microtype,booktabs,longtable,array,enumitem,xcolor}
\usepackage{fancyhdr,hyperref}
\definecolor{ink}{HTML}{18374A}
\definecolor{accent}{HTML}{236B80}
\hypersetup{colorlinks=true,linkcolor=ink,citecolor=accent,urlcolor=accent,
 pdftitle={Critical Stokes Reversion and q-Gamma Cores},
 pdfauthor={Research article prepared for Vladimir Reshetnikov}}
\pagestyle{fancy}\fancyhf{}
\fancyhead[L]{\small Critical Stokes reversion and $q$-gamma cores}
\fancyhead[R]{\small October 2026}\fancyfoot[C]{\thepage}
\setlength{\parskip}{4pt}\setlength{\parindent}{1.2em}
\setlist{itemsep=4pt,topsep=5pt}
\newcolumntype{P}[1]{>{\raggedright\arraybackslash}p{#1}}
\numberwithin{equation}{section}
\newtheorem{theorem}{Theorem}[section]
\newtheorem{proposition}[theorem]{Proposition}
\newtheorem{lemma}[theorem]{Lemma}
\newtheorem{corollary}[theorem]{Corollary}
\theoremstyle{definition}
\newtheorem{definition}[theorem]{Definition}
\newtheorem{example}[theorem]{Example}
\theoremstyle{remark}\newtheorem{remark}[theorem]{Remark}
\newcommand{\C}{\mathbb C}\newcommand{\R}{\mathbb R}
\newcommand{\N}{\mathbb N}\newcommand{\Z}{\mathbb Z}
\newcommand{\ii}{\mathrm i}\newcommand{\e}{\mathrm e}
\newcommand{\dd}{\,\mathrm d}\newcommand{\id}{\operatorname{id}}
\newcommand{\Log}{\operatorname{Log}}\newcommand{\Li}{\operatorname{Li}}
\newcommand{\Res}{\operatorname{Res}}\newcommand{\PV}{\operatorname{PV}}
\newcommand{\Borel}{\mathcal B}\newcommand{\Lap}{\mathcal L}
\newcommand{\Sp}{\mathcal S_{+}}\newcommand{\Sm}{\mathcal S_{-}}
\newcommand{\Med}{\mathcal M}\newcommand{\eps}{\varepsilon}
\newcommand{\abs}[1]{\left|#1\right|}
\newcommand{\norm}[1]{\left\lVert#1\right\rVert}
\newcommand{\coef}[2]{[#1^{#2}]}
\newcommand{\repo}{\texttt{VladimirReshetnikov/ProveIt}}
\begin{document}
\hypersetup{pageanchor=false}
\begin{titlepage}
\vspace*{10mm}
{\color{accent}\large COMPLEX ASYMPTOTICS \hfill 4 OCTOBER 2026}\par
\vspace{16mm}
{\color{ink}\fontsize{31}{36}\selectfont\bfseries Critical Stokes Reversion\\[5pt]
and $q$-Gamma Cores}\par
\vspace{9mm}
{\Large Exact Mellin--Borel completion,\\[3pt]
$q$-Pochhammer summability, and ramified inverse transseries}\par
\vspace{10mm}
{\large Prepared for Vladimir Reshetnikov}\par
\vspace{11mm}
\begin{abstract}
We develop a branch-resolved reversal calculus for sectorial complex
transseries and apply it to a moving-parameter $q$-factorial and to the
$q$-gamma function. An explicit Mellin--Borel identity reconstructs
$\log(\e^{-ah};\e^{-h})_\infty$ from its lateral Borel sums and convergent
dual products. After a rotation of variables, this proves the formulas
stated as Conjecture~5 in Fantini--Rella, arXiv:2506.08265v2. A finite
Fourier--M\"obius argument characterizes every rational-parameter weighted
sum whose median recovers the original function; the Dirichlet-character
case proves their Conjecture~1.

For inverse problems, a regular inverse transports a small jump by a
reciprocal derivative, whereas an order-$r$ critical point can change the
action from $A$ to $A/r$. We prove an all-orders analytic blow-up theorem,
a uniform fold chart, and a jet-cancellation rule. For the natural minimum
of $\log\Gamma_{\e^{-h}}(a)$, an exact correction of size
$\e^{-4\pi^2/h}$ separates inverse roots on the scale
$\e^{-2\pi^2/h}$. We compute the coefficient, the moving critical value,
and the first three ramified sectors. We also reverse the $q$-factorial
with respect to $h$, where its dual action becomes an action-$24$ sector
in the dominant inverse coordinate.

Conventional proofs and reproducible exact and high-precision checks are
provided. Stokes discontinuities, independent flat completion data,
logarithmic monodromy, and inverse-sheet branching are kept distinct.
\end{abstract}
\vfill
\noindent\small\textbf{Research status.} Theorems are proved in the specified
local domains. The new proof and applications have not been independently
refereed or formalized in a proof assistant. First-priority claims are not
made; the comparison is with the inspected sources, not an exhaustive
literature or repository audit.
\end{titlepage}
\hypersetup{pageanchor=true}\setcounter{page}{1}
\begingroup\small\setlength{\parskip}{0pt}
\tableofcontents
\endgroup
\newpage

\section{The questions answered and the source boundary}
\label{sec:scope}

\subsection{Why a new inverse chart is necessary}
The relevant difficulty is not the existence of the formal
Lagrange--B\"urmann formula. It is deciding which analytic function is
being inverted, on which sheet, and to what absolute accuracy. Three
operations that look similar in a purely symbolic calculation have
different meanings:
\[
 \text{formal reversion},\qquad
 \text{sectorial inversion},\qquad
 \text{physical-function inversion}.
\]
The first need not determine the third. A prescribed summation convention
can still miss a nonzero function flat at the expansion point. Near a
critical value this missing term can dominate the inverse error.

Our principal special-function variables are
\begin{equation}\label{eq:basic-vars}
 h=-\Log q,\qquad \Re h>0,\qquad
 A=4\pi^2,\qquad Q=\e^{-A/h},
\end{equation}
with $h\to0$ on closed subsectors of the right half-plane. The parameter
$a$ is initially real, $0<a<1$. We later continue locally to complex
neighborhoods of every positive real $a$. The moving parameter
$z=q^a=\e^{-ah}$ tends to $1$, precisely where a fixed-$z$ polylogarithmic
expansion ceases to be uniform.

\subsection{Relation to the ProveIt program}
The repository documentation was inspected at commit
\begin{center}
\texttt{b484725c893a7a9bde856151988996aef6965cb0}.
\end{center}
The canonical transseries volume and its documentation distinguish
formal inversion, dominant cores, and residual certificates; the
companion volume treats combinatorial inverses and $q$-products
\cite{repo}. The preceding report \emph{Complex Transseries Reversion at
$q$-Cusps} treats finite-action convergent units and modular Euler
quotients. Its fixed-parameter shifted-factorial inverse is explicitly
asymptotic rather than identified with a Borel sum. Moving gamma cores
and summation-compatible inversion are listed there as further questions
\cite{prior}.

We do not repeat its eta-quotient classification, cusp-data recovery, or
four-factor flatness construction. Instead, we provide a divergent
algebraic block with a fully identified analytic completion, then reverse
it in both a regular and a critical chart. The present article is
self-contained: no mathematical claim here depends on an unchecked
statement in a repository draft. The repository's scoped Lean results
do not imply formal verification of the new theorems.

\subsection{Relation to published work}
Fantini--Rella compute the rational-parameter Borel singularities and
Stokes constants of $q$-Pochhammer expansions. In the inspected version,
arXiv:2506.08265v2, dated 1 April 2026, their Conjecture~5 identifies
certain lateral sums with corrected products; their Conjecture~1 concerns
median reconstruction of character-weighted sums \cite{fr}.
Section~\ref{sec:mellin} gives a proof of that completion identity, and
Section~\ref{sec:weights} proves a stronger arbitrary-weight criterion.
The source's general modular-resurgence conjectures are \emph{not}
thereby resolved.

The Bernoulli expansions are classical; we credit McIntosh and the recent
moving-parameter treatment of Arabi Ardehali--Rosengren
\cite{mcintosh,ar}. Nonlinear stability of appropriate resurgent and
Borel-summable classes under inversion is also established mathematics
\cite{sauzin,sauzinintro}. Our proof does not present that closure theorem
as new, and does not infer analytic summability solely from formal
closure. The elementary Mellin identity below supplies the exact
completion needed for this particular family.

\subsection{Main result ledger}
\begin{enumerate}[label=\textbf{R\arabic*.},leftmargin=14mm]
\item An explicit, normally convergent meromorphic Borel kernel,
including an entire-in-$a$ representation suitable for parameter
 differentiation, and an exact lateral completion formula.
\item A proof of the two specified Fantini--Rella conjectures, with the
rotation, branch convention, and range $1\le k<N$ written explicitly.
The weighted criterion does not require multiplicativity or primitivity.
\item A local critical reversion theorem: an action $A$ becomes $A/r$
under a nondegenerate order-$r$ balance. For folds, a moving-critical-value
chart gives all orders without a singular division by the forward derivative.
\item Exact $q$-gamma completion and ramified inverse sectors near its
natural minimum, together with regular inverses in $a$ and in $h$.
\item Counterexamples to two tempting shortcuts: zero Stokes jump does
not imply zero flat correction, and arithmetic averaging does not commute
with nonlinear inversion.
\end{enumerate}

\section{Complex transseries require charts and realizations}
\label{sec:charts}

\subsection{Sectorial scales}
On a simply connected sector
\[
 S=\{0<|h|<h_0:\theta_0<\arg h<\theta_1\},
\]
choose branches of $\Log h$ and of all algebraic powers. Expressions such as
\begin{equation}\label{eq:trans-form}
 \sum_{\alpha\in\Lambda}\e^{-\alpha/h}
 h^{\beta_\alpha}(\Log h)^{m_\alpha}\widehat f_\alpha(h)
\end{equation}
become useful only after specifying the support convention and the
meaning of each coefficient block $\widehat f_\alpha$. A closed subsector
is \emph{small for an action} $\alpha$ when
$\Re(\alpha/h)>c/|h|$ there for some $c>0$.

Finite collections of small actions generate locally finite positive
projection budgets. Resonant actions must be combined. At a direction
where $\Re((\alpha-\beta)/h)=0$, two exponential magnitudes tie; that is a
\emph{dominance wall}, not by itself a Borel-summation discontinuity. A
Stokes direction requires a specified Borel singularity and a specified
summation problem. For the present tail, the positive real $h$ direction
is a Stokes direction, while $\Re(A/h)=0$ lies on the boundary of the
right half-plane.

\subsection{A formal contraction principle}
\begin{lemma}[Filtered implicit inversion]\label{lem:filtered}
Let $\mathcal A$ be a complete separated algebra over $\C$, with a
descending multiplicative filtration $I^n$. Suppose $T:I\to I$ is a
well-defined map satisfying
\[
 u-v\in I^n\quad\Longrightarrow\quad T(u)-T(v)\in I^{n+1}.
\]
Then $u=T(u)$ has exactly one solution in $I$. Picard iteration starting
at $0$ computes it to every filtration order.
\end{lemma}
\begin{proof}
The consecutive iterates differ in $I^n$ after $n$ steps, hence converge
in the complete topology. The contraction property makes $T$ continuous
and the limit fixed. Two fixed points have difference in every $I^n$,
and separatedness makes them equal.
\end{proof}
This elementary lemma is deliberately not a universal transseries-field
construction. Taylor substitution, differentiation, and action
composition must first be defined in the algebra being used.

For a regular implicit equation
$F_0(x)+R(x,\mathbf t)=y$, with $F_0'(x_0)\ne0$, the correction
$x-x_0$ has this form in the ideal generated by the flat parameters
$\mathbf t$. At a critical point the linear coefficient is not a unit.
A ramified substitution is required before applying the same principle.

\subsection{Analytic realizations and nonlinear closure}
We use two distinct analytic settings. In the first, the algebraic
blocks are actual sectorial holomorphic functions and the dependence on
independent flat parameters is convergent. Ordinary analytic implicit
inversion is sufficient. In the second, blocks are Gevrey series with
specified Borel transforms and exponential bounds. Lateral summation
then produces the actual functions to which the first setting applies.

In a class for which Borel--Laplace summation is an algebra morphism and
is compatible with analytic substitution and the implicit-function
operation, uniqueness gives
\begin{equation}\label{eq:sum-inverse}
 \mathcal S_\theta(\widehat G)
       =\bigl(\mathcal S_\theta(\widehat F)\bigr)^{-1}
\end{equation}
on matching regular inverse charts. The established closure results in
\cite{sauzin} justify this procedure under their hypotheses. Equation
\eqref{eq:sum-inverse} does not assert that the right side is the inverse
of an arbitrarily chosen special function having that expansion.
Our completion theorem determines the additional terms explicitly.

\subsection{A countable analytic version}
\begin{proposition}[Normally summable perturbations]\label{prop:countable}
Let $F_0$ be holomorphic on $|x-x_0|<R$, with $F_0'(x_0)=b\ne0$.
Suppose $d_j$ are bounded holomorphic functions on that disk and
$\sum_j |t_j|\norm{d_j}_R<\infty$. For this sum sufficiently small,
\[
 F_0(x)+\sum_j t_jd_j(x)=F_0(x_0)
\]
has one root in a fixed smaller disk, near $x_0$. That root is a
holomorphic function of the weighted $\ell^1$ parameter
$(t_j\norm{d_j}_R)_j$. Its expansion has the usual convergent homogeneous
multilinear terms.
\end{proposition}
\begin{proof}
Normal convergence gives a bounded holomorphic perturbation. Cauchy's
estimate bounds its derivative on a smaller disk by a constant times the
weighted norm. Shrink that disk until $F_0'$ is close to $b$; then shrink
the parameter ball until the map
$x\mapsto x-[F_0(x)+\sum t_jd_j(x)-F_0(x_0)]/b$
is a uniform contraction mapping the disk into itself. Iterates are
holomorphic in the Banach parameter; locally uniform convergence gives
a holomorphic fixed point. Cauchy estimates on the parameter ball give
its Taylor expansion.
\end{proof}
Thus analytic inversion can accommodate even accumulating action values
when their amplitudes are normally summable. This does not automatically
turn the resulting expression into a formally locally finite action
expansion. Analytic summability and formal support admissibility are
different requirements.

\section{Regular inversion, jump transport, and error certificates}
\label{sec:regular}

\subsection{An all-orders perturbation formula}
\begin{theorem}[Regular inverse response]\label{thm:response}
Let $f$ be univalent near $x_0$, write $g=f^{-1}$, and let $d$ be
holomorphic there. For sufficiently small independent $t$, the local
inverse $g_t$ of $f_t=f+td$ is holomorphic in $(y,t)$ and
\begin{equation}\label{eq:response}
 g_t(y)=g(y)+\sum_{n\ge1}\frac{(-t)^n}{n!}
 \frac{\mathrm d^{n-1}}{\mathrm dy^{n-1}}
 \left[g'(y)d(g(y))^n\right].
\end{equation}
The series converges normally on smaller product neighborhoods.
\end{theorem}
\begin{proof}
Put $z=f(x)$, so that $z+tD(z)=y$ with $D=d\circ g$. Classical
Lagrange inversion with the observable $g(z)$ gives
\eqref{eq:response}. Alternatively, analytic implicit inversion supplies
the convergent solution, and differentiating its defining equation
computes the same coefficients uniquely. This is the classical
observable form of Lagrange--B\"urmann, used here with an independent
flat parameter rather than a numerical perturbation of unspecified size.
\end{proof}
In source coordinates, the first terms are
\begin{equation}\label{eq:response-two}
 g_t-g=-t\frac{d}{f'}+
 t^2\left(\frac{dd'}{(f')^2}-\frac{f''d^2}{2(f')^3}\right)+O(t^3),
\end{equation}
where the right side is evaluated at $x=g(y)$. For a perturbation
$d(x,t)$ holomorphic in $t$, the analytic implicit theorem still gives
all orders; the coefficients are obtained by expanding the defining
equation, or by introducing an auxiliary multiplier before using
\eqref{eq:response}.

\subsection{Source and target Stokes maps}
Suppose two sectorial functions $F_-$ and $F_+$ are locally univalent
and have inverses $G_-$ and $G_+$ on the selected charts. There are two
useful, but different, transition conventions:
\begin{align}
 F_+&=T\circ F_- &\Longrightarrow&& G_+&=G_-\circ T^{-1},
       \label{eq:target-transition}\\
 F_+&=F_-\circ U &\Longrightarrow&& G_+&=U^{-1}\circ G_-.
       \label{eq:source-transition}
\end{align}
These identities are exact. For an additive forward jump
$F_+=F_-+\Delta$, Theorem~\ref{thm:response} gives its inverse expansion
by placing an auxiliary parameter in front of $\Delta$.

On triple overlaps, transition maps compose in their actual source or
target order. Taking inverses reverses that order. No new Stokes constant
is determined from formal coefficients alone by these identities; one
must first identify the forward sectorial functions and their jumps.

\subsection{An inverse-existence certificate}
\begin{lemma}[Disk certificate]\label{lem:certificate}
Let $f$ be holomorphic on a neighborhood of the closed disk
$D=\{|x-z|\le r\}$. Suppose $b\ne0$, $0\le\kappa<1$, and
\[
 \sup_D|1-f'(x)/b|\le\kappa,\qquad
 |f(z)-y|/|b|\le(1-\kappa)r.
\]
Then $f(x)=y$ has exactly one root in $D$, and
\begin{equation}\label{eq:certificate}
 |x-z|\le\frac{|f(z)-y|}{(1-\kappa)|b|}.
\end{equation}
\end{lemma}
\begin{proof}
The map $T(x)=x-(f(x)-y)/b$ is $\kappa$-Lipschitz on the convex disk,
by integrating its derivative along segments. The residual bound makes
$T(D)\subseteq D$. The contraction theorem gives the root and uniqueness;
$|x-z|\le\kappa|x-z|+|T(z)-z|$ gives the estimate.
\end{proof}
A bare lower bound for $|f'|$ on an arbitrary complex domain would not
justify the segment estimate. The explicit closeness to one constant
$b$ is an essential feature of this certificate.

\subsection{Averaging is not a nonlinear summation rule}
\begin{proposition}[Median and inverse do not commute]\label{prop:median}
Let $f_\pm=f\pm td/2$, and let $g_\pm=f_\pm^{-1}$. Then
\begin{equation}\label{eq:median-failure}
 \frac{g_++g_-}{2}=g+\frac{t^2}{8}
 \frac{\mathrm d}{\mathrm dy}\left[g'(y)d(g(y))^2\right]+O(t^4).
\end{equation}
Thus the arithmetic mean of inverse branches is generally not the
inverse of the arithmetic mean of the forward branches.
\end{proposition}
\begin{proof}
Substitute $t/2$ and $-t/2$ into \eqref{eq:response} and average. All
odd powers cancel; the displayed quadratic term remains.
\end{proof}
We use $\Med=(\Sp+\Sm)/2$ only for the linear logarithmic Borel problem
below. An algebra-compatible \emph{\'Ecalle median} on a full transseries
space is a different construction. It must not be replaced by arithmetic
averaging after nonlinear operations.

\section{Critical reversion and fractional actions}
\label{sec:critical-general}

\subsection{Uniform blow-up at an order-\texorpdfstring{$r$}{r} critical point}
\begin{theorem}[Critical flat-parameter reversion]\label{thm:critical}
Let $c(h)$ be holomorphic on a small punctured $h$-sector, and let
$F(x,h,t)$ be jointly holomorphic on moving neighborhoods of
$(c(h),h,0)$, with uniform $x,t$ radii and bounds. Assume
\begin{align}
 F(c+z,h,0)&=v(h)+a_r(h)z^r+O(z^{r+1}),
     &r&\ge2,\label{eq:critical-core}\\
 F(c+z,h,t)&=F(c+z,h,0)+t\,d(c+z,h)+O(t^2),
     &|a_r(h)|&\ge a_0>0.\label{eq:critical-perturb}
\end{align}
Fix compact simply connected target patches avoiding
$Y=d(c(h),h)$ uniformly. On chosen root branches, the equation
$F(x,h,t)=v(h)+tY$ has $r$ local solutions, holomorphic in
$s=t^{1/r}$, of the form
\begin{equation}\label{eq:critical-root}
 x_j=c(h)+s\,\omega_j
 \left(\frac{Y-d(c(h),h)}{a_r(h)}\right)^{1/r}+O(s^2),
 \qquad \omega_j^r=1.
\end{equation}
Each branch has a convergent expansion in $s$, locally uniformly in
$Y,h$. Truncating after $s^M$ has error $O(|s|^{M+1})$.
\end{theorem}
\begin{proof}
Set $x=c+su$, $t=s^r$, subtract $v+s^rY$, and divide by $s^r$.
The result extends holomorphically to $s=0$ and is
\[
 a_r u^r+d(c,h)-Y+s\,R(u,s,Y,h)=0.
\]
Its $r$ limiting roots are nonzero and separated on the stated compact
patches. Their $u$-derivatives, $r a_r u^{r-1}$, stay away from zero.
The analytic implicit theorem gives one branch near each limiting root,
with uniform radii after shrinking the patches. These account for all
local roots in the corresponding scaled neighborhoods, by Rouch\'e's
theorem. Cauchy's estimate in $s$ gives the remainder.
\end{proof}
For $t=\e^{-A/h}$, choose $s=\e^{-A/(rh)}$ on the sector. This proves
an action-division law rather than merely a formal square-root analogy.
It requires the nonvanishing leading balance. The excluded target value
is where another, finer critical scaling is necessary.

Comparing the unperturbed and perturbed inverse branches at the same
$y=v+tY$, on patches avoiding both $0$ and $d(c,h)$, gives
\begin{equation}\label{eq:critical-jump}
 G_{t,j}-G_{0,j}
 =s\,\omega_j a_r^{-1/r}
 \left((Y-d(c,h))^{1/r}-Y^{1/r}\right)+O(s^2).
\end{equation}
At such targets, linearizing a jump as $-\Delta/F'$ is not uniformly
valid: the nonlinear root difference is of the same order as each
root displacement.

\subsection{An exact moving-fold chart}
\begin{theorem}[Fold chart through the moving critical value]
\label{thm:morse}
Under \eqref{eq:critical-core}--\eqref{eq:critical-perturb} with $r=2$,
there are unique holomorphic functions $c_t(h)$ and $v_t(h)$ near
$c(h),v(h)$ such that $F_x(c_t,h,t)=0$ and $v_t=F(c_t,h,t)$.
There is a holomorphic local coordinate $w=w(x,h,t)$ with nonzero
$x$-derivative for which
\begin{equation}\label{eq:morse}
 F(x,h,t)=v_t(h)+w(x,h,t)^2.
\end{equation}
Consequently, on a chosen square-root cover, both inverse sheets are
\begin{equation}\label{eq:fold-exact}
 x=\chi\bigl(\pm\sqrt{y-v_t(h)},h,t\bigr),
\end{equation}
where $\chi$ is holomorphic. This representation remains valid at the
moving branch point; only the square root is ramified.
\end{theorem}
\begin{proof}
Apply the analytic implicit theorem to $F_x=0$, using $F_{xx}(c,h,0)\ne0$.
Taylor's formula factors
$F(x,h,t)-v_t=(x-c_t)^2 A(x,h,t)$ with $A$ nonvanishing near the base
point. Choose one analytic square root of $A$, and put
$w=(x-c_t)\sqrt A$. Its derivative at $c_t$ is nonzero, so it has a
holomorphic local inverse $\chi$. Uniformity follows on the same smaller
parameter neighborhoods as before.
\end{proof}
This is the one-variable parametric Morse construction, not a new
singularity-classification theorem. Its role is to provide a correctly
centered inverse-transseries chart.

Writing $F_0=F(\cdot,h,0)$, the first moving-center terms are
\begin{align}
 c_t&=c-t\frac{d'(c)}{F_0''(c)}+O(t^2),\label{eq:center-shift}\\
 v_t&=v+t d(c)+t^2\left(e(c)-\frac{d'(c)^2}{2F_0''(c)}\right)+O(t^3),
 \label{eq:value-shift}
\end{align}
when $F=F_0+t d+t^2e+O(t^3)$. All coefficients retain their full $h$
dependence.

\subsection{Three explicit ramified sectors}
At a fold, write
\begin{align*}
 F_0(c+z)&=v+\alpha_2 z^2+\alpha_3z^3+\alpha_4z^4+O(z^5),\\
 d(c+z)&=d_0+d_1z+d_2z^2+O(z^3),\qquad e(c)=e_0.
\end{align*}
For $y=v+tY$, put $s=\sqrt t$ and
$x=c+u_0s+u_1s^2+u_2s^3+O(s^4)$. Direct coefficient matching gives
\begin{align}
 u_0^2&=\frac{Y-d_0}{\alpha_2},\label{eq:u0}\\
 u_1&=-\frac{\alpha_3u_0^2+d_1}{2\alpha_2},\label{eq:u1}\\
 u_2&=-\frac{\alpha_2u_1^2+3\alpha_3u_0^2u_1+\alpha_4u_0^4
                +d_1u_1+d_2u_0^2+e_0}{2\alpha_2u_0}.
       \label{eq:u2}
\end{align}
The nonzero denominator in the last equation is exactly the simple-root
condition in the scaled chart. This recurrence continues to all orders,
and Theorem~\ref{thm:critical} proves convergence, not just consistency
of a formal calculation.

\subsection{Jet cancellation changes the action rule}
\begin{proposition}[Shared-factor cancellation]\label{prop:jet}
Suppose at the fixed target $v$ an analytic family factors exactly as
\[
 F(c+z,t)-v=z^m\bigl(z^{r-m}A(z)+tD(z,t)\bigr),
 \quad 0\le m<r,
\]
with $A(0)D(0,0)\ne0$. There are $m$ roots pinned at $c$, counted with
multiplicity, and $r-m$ moving roots with scale $t^{1/(r-m)}$ and
convergent ramified expansions.
\end{proposition}
\begin{proof}
The first factor supplies the pinned roots. Apply
Theorem~\ref{thm:critical}, or its proof, to the second factor after
$z=t^{1/(r-m)}u$.
\end{proof}
For a fold with a perturbation vanishing at $c$ but not in its first
derivative, the moving-root scale is $t$, not $\sqrt t$. The critical
value first changes at order $t^2$. Vanishing of one leading coefficient
must not be confused with the exact factorization used in this proposition.

For several commensurate actions, the lower Newton polygon of the
monomials $t^n z^m$ organizes these balances. This is the usual
Newton--Puiseux mechanism applied to a flat parameter. Arbitrary
incommensurate or accumulating action sets require additional support
and convergence hypotheses; no global Newton--Puiseux theorem for all
formal complex transseries is claimed here.

\section{The moving \texorpdfstring{$q$}{q}-factorial and its Borel transform}
\label{sec:borel}

\subsection{Forward expansion and conventions}
Define
\begin{equation}\label{eq:p-def}
 p(a,h)=\log(\e^{-ah};\e^{-h})_\infty
       =-\sum_{n\ge1}\frac{\e^{-ahn}}{n(1-\e^{-hn})}.
\end{equation}
For real $a>0$ and $\Re h>0$, the right side fixes the analytic logarithm.
The Bernoulli polynomials are defined by
\[
 \frac{s\e^{as}}{\e^s-1}=\sum_{n\ge0}B_n(a)\frac{s^n}{n!},
 \qquad B_n=B_n(0),\quad B_1=-\tfrac12.
\]
The elementary--gamma core and the formal tail are
\begin{align}
 E(a,h)&=-\frac{\pi^2}{6h}+
  \left(\frac12-a\right)\Log h+\frac12\log(2\pi)-\log\Gamma(a)
  +\frac{B_2(a)}4h,\label{eq:E}\\
 \widehat\Phi_a(h)&=\sum_{k\ge1}
 \frac{B_{2k}B_{2k+1}(a)}{2k(2k+1)!}\,h^{2k},
 \qquad \widehat p=E-\widehat\Phi_a.\label{eq:Phi}
\end{align}
This is the moving-parameter gamma expansion in the classical
$q$-factorial calculus \cite{mcintosh,ar}; we derive the exact remainder
rather than assume it.

Our Borel--Laplace normalization is
\begin{equation}\label{eq:borel-convention}
 \Borel(h^n)=\frac{\xi^{n-1}}{(n-1)!},\quad n\ge1,
 \qquad
 \Lap_\theta b(h)=\int_0^{\e^{\ii\theta}\infty}
            \e^{-\xi/h}b(\xi)\dd\xi.
\end{equation}
There is no extra $1/h$ factor. The operators $\Sp$ and $\Sm$ refer,
respectively, to contours above and below the positive real Borel poles.
The core $E$ is carried along exactly, with its chosen logarithm.

\subsection{A parameter-holomorphic kernel}
\begin{theorem}[Explicit Borel kernel]\label{thm:borel}
The Borel transform $b_a=\Borel\widehat\Phi_a$ is the germ at zero of
\begin{equation}\label{eq:borel-entire}
 b_a(\xi)=-\frac2\xi\sum_{n\ge1}
 \left\{
 \frac{\sin\bigl((a-\frac12)\xi/(2\pi n)\bigr)}
      {2\sin\bigl(\xi/(4\pi n)\bigr)}-
      \left(a-\frac12\right)\right\}.
\end{equation}
The apparent singularity at zero is removable. The sum converges
normally on compact subsets away from $\xi\in4\pi^2\Z\setminus\{0\}$,
uniformly on compact subsets of the entire $a$-plane. Its only possible
nonzero Borel singularities are simple poles at those locations.
For $N\ge1$,
\begin{equation}\label{eq:residues}
 \Res_{\xi=4\pi^2N}b_a(\xi)
 =\Res_{\xi=-4\pi^2N}b_a(\xi)
 =-\frac1\pi\sum_{d\mid N}\frac{\sin(2\pi da)}d.
\end{equation}
\end{theorem}
\begin{proof}
Set
\begin{equation}\label{eq:Tgen}
 T_a(s)=\frac{\sinh((a-\frac12)s)}{2\sinh(s/2)}-
       \left(a-\frac12\right)
       =\sum_{k\ge1}\frac{B_{2k+1}(a)}{(2k+1)!}s^{2k}.
\end{equation}
The last identity follows by taking the odd part of the Bernoulli
polynomial generating function and dividing by $s$.
Since
\[
 \frac{B_{2k}}{(2k)!}
 =\frac{2(-1)^{k+1}\zeta(2k)}{(2\pi)^{2k}},
\]
Borel transformation of \eqref{eq:Phi} gives
$-2\xi^{-1}\sum_{n\ge1}T_a(\ii\xi/(2\pi n))$, which is
\eqref{eq:borel-entire}. On compact parameter sets the summand is
$O(|\xi|^2/n^2)$ for large $n$. This proves normal convergence and
removability at zero.

At $\xi=4\pi^2N$, only the finitely many terms with $n\mid N$ have
poles. Put $d=N/n$ and use
$\sin(2\pi d(a-\frac12))=(-1)^d\sin(2\pi da)$.
The residue of the corresponding term is
$-\sin(2\pi da)/(\pi d)$. Summation gives the residue at the positive pole.
The kernel is odd in $\xi$, so residues at opposite poles are equal.
\end{proof}
The rational-parameter singularities and their divisor-sum Stokes data
are already present in \cite{fr}. Representation
\eqref{eq:borel-entire} is useful here because it remains valid when
$a$ is differentiated or moved off the real axis. A Fourier series with
coefficients $\sin(2\pi ma)$ generally diverges when $\Im a\ne0$;
we do not use such a series to claim complex-parameter uniformity.

\subsection{The real-parameter Fourier kernel and its bounds}
\begin{lemma}\label{lem:fourier-kernel}
For real $0<a<1$, one also has
\begin{equation}\label{eq:borel-fourier}
 b_a(\xi)=\frac2\pi\sum_{m,n\ge1}\frac{\sin(2\pi ma)}m
      \frac{\xi}{(4\pi^2mn)^2-\xi^2}.
\end{equation}
On each ray whose angle is bounded away from $0$ and $\pi$, this kernel
is bounded at infinity and is $O(\xi)$ at zero, uniformly for real
$a\in[0,1]$. The double series can be integrated termwise in a
convergent Laplace direction.
\end{lemma}
\begin{proof}
Use the Fourier expansion
\[
 B_{2k+1}(a)=(-1)^{k+1}
 \frac{2(2k+1)!}{(2\pi)^{2k+1}}
 \sum_{m\ge1}\frac{\sin(2\pi ma)}{m^{2k+1}},\quad k\ge1.
\]
Together with the Bernoulli-number identity in the preceding proof,
this sums a geometric series in $\xi^2$ and yields
\eqref{eq:borel-fourier}, initially near zero. Both sides continue
meromorphically, and the double series is normally convergent on compact
sets away from its poles.

If $r=|\xi|$ lies on an allowed ray, then
\[
 \left|\frac{\xi}{(4\pi^2mn)^2-\xi^2}\right|
 \le C_\theta\frac{r}{(mn)^2+r^2}.
\]
For $m\le r$ the sum over $n$ is $O(1/m)$, and for $m>r$ it is
$O(r/m^2)$. The additional factor $1/m$ makes the sums over $m$
bounded. Near zero use $\sum m^{-3}n^{-2}<\infty$. These bounds,
multiplied by the decaying Laplace exponential, justify Fubini's theorem.
\end{proof}
Consequently, the tail is Gevrey~1 and summable away from the two real
Borel rays. For complex $a$ in a fixed small neighborhood, the
parameter-holomorphic kernel gives exponential-type bounds on such
rays; sufficiently small $h$ then suffices. Integer shifts in $a$ can
also be handled by the exact $q$-factorial recurrence.

\section{An exact Mellin--Borel completion theorem}
\label{sec:mellin}

\subsection{The completion identity}
Define the convergent dual logarithms by their zero-constant-term series:
\begin{equation}\label{eq:Lpm}
 L_\pm(a,Q)=\log(\e^{\pm2\pi\ii a}Q;Q)_\infty
 =-\sum_{m\ge1}\frac{\e^{\pm2\pi\ii ma}}m\frac{Q^m}{1-Q^m}.
\end{equation}
For real $a$ they are holomorphic for $|Q|<1$. Locally in complex $a$,
the same is true whenever $|Q|\e^{2\pi|\Im a|}<1$.

\begin{theorem}[Exact lateral reconstruction]\label{thm:completion}
For $0<a<1$ real and $\Re h>0$, with the conventions above,
\begin{align}
 \Sp\widehat p(a,h)&=p(a,h)-L_-(a,Q),\label{eq:pplus}\\
 \Sm\widehat p(a,h)&=p(a,h)-L_+(a,Q).\label{eq:pminus}
\end{align}
In particular,
\begin{equation}\label{eq:median-p}
 p(a,h)=E(a,h)-\Med\widehat\Phi_a(h)
                      +\frac{L_+(a,Q)+L_-(a,Q)}2,
 \qquad \Med=\frac{\Sp+\Sm}{2},
\end{equation}
and the tail has the exact Stokes jump
\begin{align}
 \Sp\widehat\Phi_a-\Sm\widehat\Phi_a
 &=L_-(a,Q)-L_+(a,Q)\label{eq:J}\\
 &=2\ii\sum_{N\ge1}\left(\sum_{d\mid N}
                       \frac{\sin(2\pi da)}d\right)Q^N.\notag
\end{align}
These identities continue locally in $a$ near positive real parameters,
for sufficiently small $h$ on the corresponding sectorial domains.
\end{theorem}

The main point is the \emph{even} combination
$(L_++L_-)/2$, which is not determined by the discontinuity
$L_--L_+$. Computing residues alone would prove the jump but would not
prove reconstruction of the physical product. We therefore prove
\eqref{eq:median-p} independently by Mellin inversion.

\subsection{A principal-value kernel}
\begin{lemma}[Mellin transform of the symmetric kernel]\label{lem:K}
For $T>0$, let
\begin{equation}\label{eq:K}
 K(T)=\PV\int_0^\infty\e^{-Tt}\frac{t}{1-t^2}\dd t
 =\frac12\left(\e^{-T}\operatorname{Ei}(T)
                         -\e^T E_1(T)\right).
\end{equation}
Then $K(T)=O(1+|\log T|)$ at zero and $K(T)=O(T^{-2})$ at infinity.
For $0<\Re u<2$,
\begin{equation}\label{eq:K-Mellin}
 \int_0^\infty T^{u-1}K(T)\dd T
       =-\frac\pi2\Gamma(u)\cot\frac{\pi u}{2}.
\end{equation}
Consequently, for $1<c<2$,
\begin{equation}\label{eq:K-inversion}
 \frac1{2\pi\ii}\int_{c-\ii\infty}^{c+\ii\infty}
 \Gamma(u)\cot\frac{\pi u}{2}\,T^{-u}\dd u
       =-\frac2\pi K(T).
\end{equation}
\end{lemma}
\begin{proof}
The partial fraction decomposition
$t/(1-t^2)=\frac12((1-t)^{-1}-(1+t)^{-1})$
gives the exponential-integral expression, with a symmetric principal
value at $t=1$. It also gives the asserted endpoint estimates, or they
follow by splitting the integral near $0,1,\infty$.

For rigor in exchanging Mellin integration and principal value, first
exclude $(1-\delta,1+\delta)$. Away from that interval Fubini applies,
since the absolute $t$-integral of
$t^{1-\Re u}/|1-t^2|$ is finite at both endpoints when
$0<\Re u<2$. The difference between this truncated kernel and its
principal value is bounded by
$C\delta(1+T)\e^{-T/2}$ for $0<\delta<1/2$, after subtracting the
pole numerator's value at $t=1$. This bound is integrable against
$T^{\Re u-1}\dd T$. Hence the limit commutes with Mellin integration.
The transform is
\[
 \Gamma(u)\PV\int_0^\infty\frac{t^{1-u}}{1-t^2}\dd t
 =\frac{\Gamma(u)}2\PV\int_0^\infty\frac{v^{-u/2}}{1-v}\dd v
 =-\frac\pi2\Gamma(u)\cot\frac{\pi u}{2}.
\]
The last principal-value identity follows, for example, by splitting at
$1$, substituting $v\mapsto1/v$ in the second piece, and applying the
partial fraction expansion of the cotangent. Mellin inversion is valid
in this strip; its right side has exponential decay on vertical lines
from $\Gamma(u)$, while the cotangent is bounded there. This proves
\eqref{eq:K-inversion}.
\end{proof}

\subsection{Mellin representation of the product}
\begin{lemma}\label{lem:MB}
For $h>0$, $a>0$, and $c_0>1$,
\begin{equation}\label{eq:MB-initial}
 p(a,h)=-\frac1{2\pi\ii}
 \int_{c_0-\ii\infty}^{c_0+\ii\infty}
       \Gamma(s)\zeta(s+1)\zeta(s,a)h^{-s}\dd s.
\end{equation}
For $0<a<1$, shift the contour to $-2<\Re s<-1$. The residues crossed
are exactly $E(a,h)$, so that
\begin{equation}\label{eq:MB-remainder}
 p(a,h)-E(a,h)=-\frac1{2\pi\ii}
 \int_{-c-\ii\infty}^{-c+\ii\infty}
       \Gamma(s)\zeta(s+1)\zeta(s,a)h^{-s}\dd s,
 \quad1<c<2.
\end{equation}
\end{lemma}
\begin{proof}
Expand the logarithm as a double sum
$-\sum_{n\ge1,j\ge0}\e^{-hn(j+a)}/n$ and insert Mellin inversion of
the exponential. Absolute convergence for $\Re s>1$ yields
\eqref{eq:MB-initial}. Standard vertical-strip bounds for the zeta
functions and exponential decay of $\Gamma$ justify the contour shift.
The pole at $s=1$ contributes $-\pi^2/(6h)$.
At $s=0$, use
\[
 \zeta(0,a)=\frac12-a,\qquad
 \zeta'(0,a)=\log\Gamma(a)-\frac12\log(2\pi),
\]
and $\Gamma(s)\zeta(1+s)=s^{-2}+O(1)$.
The residue, with the outer minus sign, is
$(\frac12-a)\log h-\log\Gamma(a)+\frac12\log(2\pi)$.
At $s=-1$, the residue is $B_2(a)h/4$ after the outer minus sign.
These standard zeta identities are recorded in \cite{dlmfzeta}.
\end{proof}

\subsection{Functional equations split the remainder}
For real $0<a<1$, set
\[
 C_u(a)=\sum_{m\ge1}\frac{\cos(2\pi ma)}{m^u},\qquad
 S_u(a)=\sum_{m\ge1}\frac{\sin(2\pi ma)}{m^u}.
\]
On the line in \eqref{eq:MB-remainder}, the Hurwitz Fourier formula and
the Riemann functional equation give
\begin{align}
 \zeta(s,a)&=2\Gamma(1-s)(2\pi)^{s-1}
 \left(\sin\frac{\pi s}{2}\,C_{1-s}(a)
             +\cos\frac{\pi s}{2}\,S_{1-s}(a)\right),\label{eq:hurwitzFE}\\
 \zeta(s+1)&=2^{s+1}\pi^s\cos\frac{\pi s}{2}
                          \Gamma(-s)\zeta(-s).\label{eq:riemannFE}
\end{align}
The first is DLMF~25.11.9 in changed variables \cite{dlmfzeta}; the
second is the standard functional equation \cite{dlmfriemann}.
Using $\Gamma(s)\Gamma(1-s)=\pi/\sin(\pi s)$ simplifies their product
to the particularly useful identity
\begin{equation}\label{eq:FE-product}
 \Gamma(s)\zeta(s+1)\zeta(s,a)
 =(4\pi^2)^s\Gamma(-s)\zeta(-s)
 \left(C_{1-s}(a)+\cot\frac{\pi s}{2}\,S_{1-s}(a)\right).
\end{equation}
Put $X=4\pi^2/h$ and $u=-s$ in \eqref{eq:MB-remainder}. Then
\begin{equation}\label{eq:MB-dual}
 p-E=-\frac1{2\pi\ii}\int_{(c)}
 \Gamma(u)\zeta(u)
 \left(C_{1+u}(a)-\cot\frac{\pi u}{2}\,S_{1+u}(a)\right)
 X^{-u}\dd u,\quad1<c<2.
\end{equation}
All Dirichlet series here converge absolutely, and all vertical
integrals converge absolutely. We may therefore integrate the two
double series separately. Mellin inversion of $\Gamma(u)$ evaluates
the cosine part as
\begin{equation}\label{eq:cos-completion}
 -\sum_{m,n\ge1}\frac{\cos(2\pi ma)}m\e^{-Xmn}
                    =\frac{L_+(a,Q)+L_-(a,Q)}2.
\end{equation}
Lemma~\ref{lem:K} evaluates the sine part as
\begin{equation}\label{eq:sine-completion}
 -\frac2\pi\sum_{m,n\ge1}\frac{\sin(2\pi ma)}m K(Xmn).
\end{equation}
The series in \eqref{eq:sine-completion} converges absolutely: for fixed
$X>0$, its tail is bounded by a constant times
$\sum m^{-3}n^{-2}$.

\subsection{Identification with the lateral Borel sums}
For a single pole pair, contour deformation above or below the positive
pole gives
\begin{equation}\label{eq:single-pole-laplace}
 \int_{0}^{\e^{\pm\ii\theta}\infty}
 \e^{-\xi/h}\frac{\xi}{A_0^2-\xi^2}\dd\xi
   =K(A_0/h)\ \pm\ \frac{\pi\ii}{2}\e^{-A_0/h},
 \quad h,A_0>0,
\end{equation}
where $0<\theta<\pi/2$, and deformation is within the indicated
half-plane. The positive-pole residue is $-1/2$.
By Lemma~\ref{lem:fourier-kernel}, we can sum the Laplace integrals
termwise. The resulting series of $K$ terms and residue exponentials
are absolutely convergent. Therefore
\begin{equation}\label{eq:median-kernel}
 \Med\widehat\Phi_a(h)
 =\frac2\pi\sum_{m,n\ge1}\frac{\sin(2\pi ma)}m K(Xmn).
\end{equation}
Equations \eqref{eq:MB-dual}--\eqref{eq:median-kernel} prove
\eqref{eq:median-p} for positive $h$.

The difference of the two contours is clockwise around each positive
pole, so its residue contribution is $-2\pi\ii$ times the residue.
Together with \eqref{eq:residues}, this proves \eqref{eq:J}, including
its sign. Substitution into the median identity gives
\eqref{eq:pplus}--\eqref{eq:pminus}.

For completeness, the lateral sums define holomorphic functions for
$\Re h>0$: choose a small positive, respectively negative, Borel ray
whose Laplace exponential decays, and deform it within its half-plane.
On overlapping $h$-domains the contours give the same analytic function.
Lemma~\ref{lem:fourier-kernel} supplies the required bounds. The
right-hand sides of \eqref{eq:pplus}--\eqref{eq:pminus} are also
holomorphic there. The identity theorem extends the equality from
$h>0$ to that entire half-plane.

Finally, \eqref{eq:borel-entire} is normally convergent in complex $a$.
Local exponential bounds allow differentiation under the Laplace
integral for small $h$. Analytic continuation in $a$, followed when
necessary by
\begin{equation}\label{eq:p-recurrence}
 p(a+1,h)=p(a,h)-\log(1-\e^{-ah}),
\end{equation}
extends the identities to the stated positive-parameter neighborhoods.
This finishes the proof of Theorem~\ref{thm:completion}.
\hfill$\square$

\section{The conjecture comparison and a complete weight criterion}
\label{sec:weights}

\subsection{Exact translation of Conjecture 5}
In the notation of \cite{fr}, let $N\ge2$, $1\le k<N$, and
\[
 g_{k,N}(y)=\log(q^k;q^N)_\infty,\qquad
 f_{k,N}(y)=\log(\e^{2\pi\ii k/N};q)_\infty,
 \qquad q=\e^{2\pi\ii y},\quad \Im y>0.
\]
Set
\begin{equation}\label{eq:FR-rotation}
 a=k/N,\qquad h=-2\pi\ii Ny,\qquad
 Q=\e^{-4\pi^2/h}=\e^{-2\pi\ii/(Ny)}.
\end{equation}
Then $g_{k,N}(y)=p(a,h)$ and
\begin{align*}
 L_+(a,Q)&=f_{k,N}(-1/(Ny))-\log(1-\e^{2\pi\ii k/N}),\\
 L_-(a,Q)&=f_{N-k,N}(-1/(Ny))-\log(1-\e^{-2\pi\ii k/N}).
\end{align*}
Here the subtracted constants mean the analytically determined product
factor at dual index zero. Writing the difference as $L_\pm$ avoids an
uncontrolled multiple of $2\pi\ii$ from separately chosen logarithms.

The positive real Borel direction in the $y$ variable rotates to the
negative imaginary direction in the $h$ variable, hence to our lower
lateral continuation. The negative real $y$ direction rotates to the
upper one. Theorem~\ref{thm:completion} therefore gives
\begin{align}
 s_0(\widetilde g_{k,N})(y)
 &=g_{k,N}(y)-f_{k,N}(-1/(Ny))
       +\log(1-\e^{2\pi\ii k/N}), &&\Re y>0,\label{eq:FR60a}\\
 s_\pi(\widetilde g_{k,N})(y)
 &=g_{k,N}(y)-f_{N-k,N}(-1/(Ny))
       +\log(1-\e^{-2\pi\ii k/N}), &&\Re y<0.\label{eq:FR60b}
\end{align}
Both domains are intersected with $\Im y>0$. These are precisely
(60a)--(60b), labeled Conjecture~5 in arXiv:2506.08265v2
\cite{fr}. Thus the Mellin calculation supplies the missing analytic
identification in that statement. The excluded index $k=0$ modulo $N$
would make the isolated symbol $f_{0,N}$ singular; it must not be inserted
into \eqref{eq:FR60a} as a finite logarithm.

\subsection{Beyond Dirichlet characters}
\begin{theorem}[Antisymmetry is necessary and sufficient]
\label{thm:weights}
Fix $N\ge2$ and arbitrary complex weights $w_1,\ldots,w_{N-1}$.
Put $w_0=0$ and
\[
 P_w(h)=\sum_{k=1}^{N-1}w_kp(k/N,h),\qquad
 \widehat P_w=\sum_{k=1}^{N-1}w_k\widehat p(k/N,h).
\]
Then
\begin{equation}\label{eq:weight-iff}
 P_w=\Med\widehat P_w\text{ identically}
 \quad\Longleftrightarrow\quad
 w_k=-w_{N-k}\quad(1\le k<N).
\end{equation}
If the equality fails, its first nonzero dual coefficient occurs at some
power $Q^n$ with $1\le n\le N$.
\end{theorem}
\begin{proof}
By Theorem~\ref{thm:completion}, the defect is
\begin{equation}\label{eq:weighted-defect}
 C_w(Q)=P_w-\Med\widehat P_w
  =-\sum_{m\ge1}\frac{c_w(m)}m\frac{Q^m}{1-Q^m},
 \quad c_w(m)=\sum_{k=1}^{N-1}w_k\cos(2\pi km/N).
\end{equation}
Antisymmetry of the weights makes each $c_w(m)$ zero.
Conversely, if $C_w=0$, its coefficient at $Q^n$ gives
$\sum_{d\mid n}c_w(d)/d=0$. M\"obius inversion yields
\begin{equation}\label{eq:mobius-c}
 c_w(n)=-n\sum_{d\mid n}\mu(n/d)[Q^d]C_w(Q)=0.
\end{equation}
The sequence $c_w(n)$ is the discrete Fourier transform of the even
part $(w_k+w_{-k})/2$ of the cyclic weight vector. Invertibility of the
finite Fourier transform implies that this even part is zero.

If it is nonzero, some $c_w(n)$ with $1\le n\le N$ is nonzero by
periodicity. Choose the least such positive $n$. All proper-divisor
terms in the coefficient at $Q^n$ vanish, leaving $-c_w(n)/n\ne0$.
This proves the finite detection bound.
\end{proof}
A nonzero Dirichlet character satisfies $\chi(-k)=\chi(-1)\chi(k)$.
For $N\ge2$ its value at the zero residue is zero. Hence
\eqref{eq:weight-iff} says exactly that reconstruction holds if and only
if the character is odd. Together with \eqref{eq:FR-rotation}, this
proves the statement labeled Conjecture~1 in \cite{fr}. Primitivity is
not needed for this summability criterion.

\subsection{Examples and what the theorem does not say}
For $N=3$, weights $(1,-1)$ give exact median reconstruction, while
$(1,1)$ give a nonzero even defect. For $N=4$, antisymmetry requires
$w_1=-w_3$ and $w_2=0$. More generally the admissible weight space has
dimension $\lfloor(N-1)/2\rfloor$ over $\C$.

The result classifies reconstruction of these \emph{linear logarithmic}
combinations. Exponentiating an arithmetic mean or averaging inverses
would be a different operation. Theorem~\ref{thm:weights} neither
constructs an algebra-compatible median for all transseries nor settles
the broader modular-resurgence conjectures in \cite{fr}.

\section{An exact \texorpdfstring{$q$}{q}-gamma completion and its moving gamma core}
\label{sec:qgamma}

\subsection{The regularized observable}
For $q=\e^{-h}$, define the standard $q$-gamma function \cite{dlmfqgamma}
by
\begin{equation}\label{eq:qgamma}
 \Gamma_q(a)=\frac{(q;q)_\infty(1-q)^{1-a}}{(q^a;q)_\infty},
 \qquad F(a,h)=\log\Gamma_{\e^{-h}}(a).
\end{equation}
For positive real $a,h$ the logarithm is real, and local complex
continuations use that branch. Cancellation of the essential and
logarithmic terms in \eqref{eq:E} gives the formal expansion
\begin{align}
 \widehat F(a,h)&=\log\Gamma(a)+E_\Gamma(a,h)+\widehat\Phi_a(h),
                  \label{eq:Fhat}\\
 E_\Gamma(a,h)&=(1-a)\Log\frac{1-\e^{-h}}h
                  +\frac{B_2(1)-B_2(a)}4h.\label{eq:Egamma}
\end{align}
The function $E_\Gamma$ is convergent near $h=0$; the generally divergent
part is entirely in $\widehat\Phi_a$.

\begin{theorem}[Exact gamma completion]\label{thm:gamma-completion}
Let $B_\pm(a,h)=\mathcal S_\pm\widehat F(a,h)$ and
$H(a,h)=(B_+(a,h)+B_-(a,h))/2$. In the local domains just specified,
\begin{align}
 F&=B_++\log(Q;Q)_\infty-L_-(a,Q)\label{eq:gamma-plus}\\
  &=B_-+\log(Q;Q)_\infty-L_+(a,Q),\label{eq:gamma-minus}\\
 F&=H+D(a,Q),\label{eq:gamma-median}
\end{align}
where
\begin{align}
 D(a,Q)&=\log(Q;Q)_\infty-\frac{L_+(a,Q)+L_-(a,Q)}2\notag\\
 &=\sum_{m\ge1}\frac{\cos(2\pi ma)-1}{m}\frac{Q^m}{1-Q^m}
   =\sum_{n\ge1}\delta_n(a)Q^n,\label{eq:D}\\
 \delta_n(a)&=\sum_{d\mid n}\frac{\cos(2\pi da)-1}{d}.
                       \label{eq:delta-n}
\end{align}
The lateral jump is
\begin{equation}\label{eq:gamma-J}
 B_+-B_-=J(a,Q)=L_-(a,Q)-L_+(a,Q).
\end{equation}
\end{theorem}
\begin{proof}
Take the limit $a\to1$ in Theorem~\ref{thm:completion}.
The tail vanishes there, because $B_{2k+1}(1)=0$ for $k\ge1$, and
$L_\pm(1,Q)=\log(Q;Q)_\infty$. This gives
$p(1,h)=E(1,h)+\log(Q;Q)_\infty$.
Insert this and \eqref{eq:pplus}--\eqref{eq:pminus} into
$F=p(1,h)+(1-a)\log(1-\e^{-h})-p(a,h)$.
The core difference is \eqref{eq:Egamma}; averaging and subtracting
give the remaining identities.

For parameters above $1$, use
\begin{equation}\label{eq:gamma-recurrence}
 F(a+1,h)-F(a,h)=\Log\frac{1-\e^{-ah}}{1-\e^{-h}}.
\end{equation}
Its nonconstant small-$h$ expansion is convergent locally. Both $D$ and
$J$ are periodic in $a$, so the completion identities persist under the
shift. Complex neighborhoods follow by analytic continuation and the
parameter-holomorphic Borel kernel.
\end{proof}

\subsection{A completely invisible nonperturbative term}
At $a=1/2$, all the Bernoulli polynomials $B_{2k+1}(a)$ in
\eqref{eq:Phi} vanish. Thus the algebraic tail is zero and so is its
Stokes jump. Nevertheless,
\begin{equation}\label{eq:half-invisible}
 D(1/2,Q)=\log(Q;Q)_\infty-\log(-Q;Q)_\infty
          =-2Q-2Q^2+O(Q^3)\ne0.
\end{equation}
Here $\widehat F$ is convergent near $h=0$, but its sum is not the
physical $F$: an essential flat term remains. This explicit example
shows that neither an all-orders algebraic expansion nor its complete
Stokes discontinuity determines the original function without additional
normalization data.

For a positive integer $a$, in contrast, $D=J=0$ and
$\Gamma_q(a)=\prod_{j=1}^{a-1}(1-q^j)/(1-q)$ is a finite product.
The difference between half-integers and integers is therefore visible
in the flat completion even when both have zero Stokes jump.

\subsection{Algebraic coefficients and parameter uniformity}
Expanding \eqref{eq:Fhat} gives
\begin{equation}\label{eq:F-two}
 F(a,h)=\log\Gamma(a)+f_1(a)h+f_2(a)h^2+O(h^4),
\end{equation}
with
\begin{equation}\label{eq:f1f2}
 f_1(a)=-\frac{(a-1)(a-2)}4,
 \qquad f_2(a)=\frac{2a^3-3a^2-5a+6}{144}.
\end{equation}
There is no cubic term; after the linear term the formal expansion has
only even powers. Formula \eqref{eq:F-two} is uniform on compact complex
neighborhoods of a positive real parameter, on smaller $h$-sectors.
The same holds after any fixed number of $a$-derivatives. One proof
uses the parameter-holomorphic Borel kernel, a uniform Laplace estimate,
and Cauchy's estimate in $a$. The exact completion $D$ is exponentially
small on these compact sets.

At $a=3$, the finite formula $\Gamma_q(3)=1+q$ gives
$\log(1+\e^{-h})=\log2-h/2+h^2/8+O(h^4)$, agreeing with
\eqref{eq:f1f2}. Such checks audit the normalization but do not replace
the Mellin proof.

\subsection{Explicit flat-tail bounds}
If $|\Im a|\le M$ and $\rho=|Q|\e^{2\pi M}<1$, then
\begin{equation}\label{eq:D-bound}
 |D(a,Q)|\le\frac{-2\log(1-\rho)}{1-|Q|}.
\end{equation}
Indeed, bound $|\cos(2\pi ma)-1|$ by $2\e^{2\pi mM}$ in
\eqref{eq:D}. An action-coefficient bound is
$|\delta_n(a)|\le2(n+1)\e^{2\pi nM}$, so truncating after $Q^K$ gives
\begin{equation}\label{eq:D-tail}
 \left|D-\sum_{n=1}^{K}\delta_n(a)Q^n\right|
 \le2\rho^{K+1}\left(\frac{K+2}{1-\rho}
                         +\frac{\rho}{(1-\rho)^2}\right).
\end{equation}
The same type of estimate applies to $J$, and parameter derivatives
follow by Cauchy's estimate on a slightly larger $a$-neighborhood.
These bounds can be inserted directly into the inverse certificates;
the numerical program is not needed to justify them.

\section{Regular inverses in the gamma parameter and in \texorpdfstring{$h$}{h}}
\label{sec:applications-regular}

\subsection{Reversing the gamma parameter}
Let $a_0$ be a positive real number with $\psi(a_0)\ne0$, and fix
$y=\log\Gamma(a_0)$ on a local complex target neighborhood. The inverse
of $F(a,h)=y$ has the all-orders algebraic expansion
\begin{align}
 a(y,h)={}&a_0-\frac{f_1(a_0)}{\psi(a_0)}h\notag\\
 &+\left(-\frac{f_2(a_0)}{\psi(a_0)}
       +\frac{f_1(a_0)f_1'(a_0)}{\psi(a_0)^2}
       -\frac{\psi_1(a_0)f_1(a_0)^2}{2\psi(a_0)^3}\right)h^2
       +O(h^3),\label{eq:a-inverse}
\end{align}
where $\psi_j$ denotes the order-$j$ polygamma function, so
$\psi_1=\psi'$ and $\psi_2=\psi''$. Coefficients at every order follow
uniquely by substituting into \eqref{eq:Fhat}. Uniform forward
remainders and Lemma~\ref{lem:certificate} show that this formal series
is the asymptotic expansion of the analytic local inverse. In a chosen
summable algebra, regular implicit closure supplies the corresponding
lateral inverse \cite{sauzin}; the physical correction must still be added.

For that correction, let $a_H(y,h)$ be the actual inverse of the
median-forward function $H$, and work on a compact regular chart where
$|H_a(a_H,h)|$ stays positive. Theorem~\ref{thm:gamma-completion} and
analytic inversion in the independent variable $Q$ give
\begin{equation}\label{eq:a-flat}
 a_F-a_H=
 \frac{1-\cos(2\pi a_H)}{H_a(a_H,h)}Q+O(Q^2).
\end{equation}
The coefficient is an actual function of $h$. Replacing it by its
leading algebraic value generally changes the error to $O(hQ)$,
which is much larger than $O(Q^2)$. This is why full coefficient blocks
are retained in an exponentially accurate inverse.

Similarly, if $G_\pm$ are regular inverses of $B_\pm$, then
\begin{equation}\label{eq:regular-Stokes}
 G_+-G_-=-\frac{2\ii\sin(2\pi G_-)}{(B_-)_a(G_-,h)}Q+O(Q^2),
\end{equation}
and Theorem~\ref{thm:response} supplies every higher jump term. The
physical correction in \eqref{eq:a-flat} is not the same as the Stokes
jump in \eqref{eq:regular-Stokes}.

\subsection{Reversing the moving factorial with respect to \texorpdfstring{$h$}{h}}
Fix $a\in(0,1)$ and put
\begin{equation}\label{eq:h-core-vars}
 A_0=\pi^2/6,\quad \beta=\tfrac12-a,\quad
 C(a)=\frac{\sqrt{2\pi}}{\Gamma(a)},\quad x=A_0/h.
\end{equation}
Given a small nonzero value $P=\exp p(a,h)$, choose the logarithmic
sheet
\[
 \ell=\Log(C(a)/P)+\beta\Log A_0.
\]
The physical inverse equation is
\begin{equation}\label{eq:h-inverse-phase}
 \ell=x+\beta\Log x-\mathcal R(a,A_0/x)
                     -C_q(a,\e^{-24x}),
\end{equation}
where
\[
 \mathcal R(a,h)=\frac{B_2(a)}4h-\Med\widehat\Phi_a(h),
 \qquad C_q(a,Q)=\frac{L_+(a,Q)+L_-(a,Q)}2.
\]
The exponent $24$ is the exact action ratio $4\pi^2/A_0$.
This is a logarithmic inverse chart on a cover, not a single-valued
inverse on the punctured $P$-plane.

The elementary core $x+\beta\Log x=\ell$ has the selected solution
$x\sim\ell$. For $\beta\ne0$, it can be written
\begin{equation}\label{eq:lambert-core}
 x_0=\beta W\left(\beta^{-1}\e^{\ell/\beta}\right),
\end{equation}
where the Lambert branch is the one satisfying that asymptotic relation.
For real large $\ell$, it is $W_0$ when $\beta>0$ and $W_{-1}$ when
$\beta<0$. Complex continuation follows the chosen inverse chart rather
than a global fixed branch label. At $\beta=0$, $x_0=\ell$.

Set
\[
 r_1=\frac{A_0 B_2(a)}4,\qquad
 r_2=-\frac{A_0^2 B_3(a)}{72},\qquad L=\Log\ell.
\]
Formal reversion of \eqref{eq:h-inverse-phase} through the algebraic
orders gives
\begin{align}
 x\sim{}&\ell-\beta L+\frac{\beta^2L+r_1}{\ell}\notag\\
 &+\frac{\tfrac12\beta^3L^2-\beta^3L+\beta r_1L-\beta r_1+r_2}
               {\ell^2}+\cdots,\label{eq:x-PL}\\
 h\sim{}&\frac{A_0}{\ell}+\frac{A_0\beta L}{\ell^2}
       +\frac{A_0(\beta^2L^2-\beta^2L-r_1)}{\ell^3}+\cdots.
                \label{eq:h-PL}
\end{align}
Substituting successively determines every polynomial in $\Log\ell$.
The coefficient of the next unknown is one in the normalized $x$
equation. Finite-order sectorial remainders and the disk certificate
prove the corresponding asymptotic inverse statements. These expansions
are generally not convergent in the algebraic index.

\subsection{A convergent lift for the complete exponential sectors}
\begin{proposition}[Action-$24$ inverse lift]\label{prop:24}
Let
\[
 f_M(x)=x+\beta\Log x-\mathcal R(a,A_0/x),
 \qquad f_M(X)=\ell,
\]
with $X$ on the selected large regular chart. There is a function
$u=u(X,t)$ holomorphic for $|t|<t_0$, with $u(X,0)=0$, such that the
physical inverse is
\begin{equation}\label{eq:24-exact}
 x=X+u(X,\e^{-24X}),\qquad
 h=\frac{A_0}{X+u(X,\e^{-24X})}.
\end{equation}
The radius $t_0$ may be chosen uniformly for large $X$ on a closed
subsector. Moreover,
\begin{align}
 u(X,t)&=-\frac{\cos(2\pi a)}{f_M'(X)}t+O(t^2),\label{eq:24-u}\\
 h-\frac{A_0}{X}&=
  \frac{A_0\cos(2\pi a)}{X^2f_M'(X)}\e^{-24X}
          +O\left(X^{-2}\e^{-48X}\right).
          \label{eq:24-h}
\end{align}
All exponential-sector coefficients are obtainable by analytic implicit
reversion, with geometric tails in $t$.
\end{proposition}
\begin{proof}
Substitute $x=X+u$ into \eqref{eq:h-inverse-phase}. Its equation is
\begin{equation}\label{eq:24-lift}
 f_M(X+u)-f_M(X)-C_q(a,t\e^{-24u})=0.
\end{equation}
For $u$ in a fixed small disk, this is holomorphic in $u,t$.
At $(u,t)=(0,0)$ its $u$-derivative is
$f_M'(X)=1+O(X^{-1})$, uniformly on the chosen subsector.
All higher local bounds are uniform there, so the analytic implicit
construction has a uniform $t$-radius. The coefficient of $Q$ in
$C_q$ is $-\cos(2\pi a)$, proving \eqref{eq:24-u}.
Expanding $A_0/(X+u)$ proves \eqref{eq:24-h}. Cauchy's estimate yields
geometric remainders for the convergent lifted series.
\end{proof}
The factor $\e^{-24u}$ in \eqref{eq:24-lift} is important: the
exponential action is evaluated at the \emph{unknown} inverse point.
Freezing it at $X$ beyond first order would miss action-feedback terms.
The coordinate $x=A_0/h$ keeps this feedback uniformly analytic even
though differentiation of $\e^{-A/h}$ in $h$ produces powers of $h^{-2}$.

At $a=1/4$ or $3/4$, the action-$24$ coefficient vanishes. The first
coefficient of $C_q$ is then $Q^2/2$, so the first physical inverse
correction has action $48$ in $X$ and equals
$-A_0\e^{-48X}/(2X^2 f_M'(X))$ to leading sector order. This is an
amplitude cancellation, not a disappearance of the dual scale itself.

\section{The natural \texorpdfstring{$q$}{q}-gamma minimum and half-action inverse roots}
\label{sec:natural-fold}

\subsection{Existence and the algebraic moving center}
For real $h>0$, differentiation of the convergent logarithmic product
gives
\begin{align}
 F_a(a,h)&=-\log(1-\e^{-h})
          -h\sum_{n\ge0}\frac1{\e^{h(n+a)}-1},\label{eq:Fprime}\\
 F_{aa}(a,h)&=h^2\sum_{n\ge0}
   \frac{\e^{h(n+a)}}{(\e^{h(n+a)}-1)^2}>0,\qquad a>0.
                         \label{eq:Fsecond}
\end{align}
The first derivative tends to $-\infty$ at $a=0$ and to
$-\log(1-\e^{-h})>0$ at infinity. Thus the physical function has
exactly one positive critical point, a strict global minimum.

Let $a_*$ be the unique positive zero of $\psi(a)$. Then
\begin{equation}\label{eq:astar}
 1<a_*<\frac32,\qquad
 a_*=1.4616321449683623412626595423\ldots.
\end{equation}
Uniqueness follows from $\psi_1(a)>0$; the endpoint signs follow from
$\psi(1)=-\gamma$ and $\psi(3/2)=2-\gamma-2\log2$ \cite{dlmfgamma}.
The elementary bounds $\gamma<3/5$ and $\log2<7/10$ make the latter
strictly positive. The accompanying exact checks verify these bounds
using finite rational sums; see Section~\ref{sec:verification}.
In particular, neither $\sin(2\pi a_*)$ nor $1-\cos(2\pi a_*)$ vanishes.

For small $h$ on a sector, let $c(h)$ denote the critical point of the
\emph{median-forward} function $H$ near $a_*$, and put
\begin{equation}\label{eq:median-center}
 v(h)=H(c(h),h),\qquad \alpha(h)=\frac12 H_{aa}(c(h),h).
\end{equation}
Uniform parameter asymptotics and the analytic implicit theorem give
this center, with $\alpha(h)\to\psi_1(a_*)/2\ne0$.
It is not defined as an arithmetic median of critical points of $B_\pm$.
That different operation would encounter Proposition~\ref{prop:median}.

Writing $c=a_*+c_1h+c_2h^2+O(h^3)$, we obtain
\begin{align}
 c_1&=\frac{2a_*-3}{4\psi_1(a_*)},\label{eq:c1}\\
 c_2&=-\frac{\tfrac12\psi_2(a_*)c_1^2-\tfrac12c_1+f_2'(a_*)}
                 {\psi_1(a_*)},\label{eq:c2}\\
 v(h)&=\log\Gamma(a_*)+f_1(a_*)h
    +\left(f_2(a_*)-\frac{f_1'(a_*)^2}{2\psi_1(a_*)}\right)h^2
       +O(h^3).\label{eq:vexp}
\end{align}
These formulas follow by expanding $H_a(c,h)=0$ and $H(c,h)$.
The same recursive calculation gives all algebraic orders. The full
functions $c(h),v(h),\alpha(h)$, rather than finite truncations of them,
are used in the exponential statements below.

\subsection{The physical critical value is different}
Write $c_F(h)$ and $v_F(h)$ for the actual positive critical point and
value of $F$. Since
$F=H+\sum_{n\ge1}\delta_n(a)Q^n$, equations
\eqref{eq:center-shift}--\eqref{eq:value-shift} give
\begin{align}
 c_F&=c-\frac{\delta_1'(c)}{H_{aa}(c,h)}Q+O(Q^2)
     =c+\frac{2\pi\sin(2\pi c)}{H_{aa}(c,h)}Q+O(Q^2),
                     \label{eq:qcriticalpoint}\\
 v_F&=v+\delta_1(c)Q+
 \left(\delta_2(c)-\frac{\delta_1'(c)^2}{2H_{aa}(c,h)}\right)Q^2
      +O(Q^3).\label{eq:qcriticalvalue}
\end{align}
Because $\delta_1(c)=\cos(2\pi c)-1<0$ for small positive $h$,
the physical minimum lies below the median-forward one. The error in
its value is exponentially small, but it is the leading datum for the
inverse in this target window.

\begin{theorem}[Two real inverse roots with half the forward action]
\label{thm:qfold}
For all sufficiently small positive $h$, the equation
\begin{equation}\label{eq:physical-level}
 F(a,h)=v(h)
\end{equation}
has exactly two positive real roots. They have the convergent lifted
expansions
\begin{equation}\label{eq:qfold-leading}
 a_\pm(h)=c(h)\ \pm\
 \sqrt{\frac{1-\cos(2\pi c(h))}{\alpha(h)}}\,
     \e^{-2\pi^2/h}+O(\e^{-4\pi^2/h}).
\end{equation}
The coefficient tends to
\begin{equation}\label{eq:Cstar}
 C_*=
 \sqrt{\frac{2(1-\cos(2\pi a_*))}{\psi_1(a_*)}}
 =2.0183815070093546945009855024\ldots.
\end{equation}
More generally, for $y=v(h)+QY$, every branch has an all-orders
expansion in $Q^{1/2}$ on compact complex $Y$-patches avoiding
$Y=\delta_1(c(h))$ uniformly.
\end{theorem}
\begin{proof}
Use the independent-parameter family
$\mathcal F(a,h,t)=H(a,h)+D(a,t)$. It is uniformly holomorphic near
$(a,t)=(c(h),0)$, and $\mathcal F_{aa}(c,h,0)=2\alpha(h)\ne0$.
Theorem~\ref{thm:critical} with $r=2$, $t=Q$, and $Y=0$ yields
\eqref{eq:qfold-leading}. For positive $h$, the square root has a
positive real argument, so both local branches are real and distinct.
Furthermore, $F(c,h)=v+D(c,Q)<v$: every real term in
\eqref{eq:D} is nonpositive, with at least one strictly negative.
Strict convexity \eqref{eq:Fsecond}, together with divergence of $F$
at both positive endpoints, shows that these are exactly the two
positive roots. The limiting coefficient follows from
$c\to a_*$ and $2\alpha\to\psi_1(a_*)$. The all-orders conclusion is
the convergent implicit expansion of Theorem~\ref{thm:critical}.
\end{proof}
The coefficient in \eqref{eq:qfold-leading} may be replaced by $C_*$
only at the cost of an $O(h\e^{-2\pi^2/h})$ error. That error is much
larger than $O(\e^{-4\pi^2/h})$; the stronger displayed remainder
requires retaining the moving coefficient.

\subsection{The first three sectors explicitly}
In \eqref{eq:u0}--\eqref{eq:u2}, substitute
\[
 \alpha_j=\frac{\partial_a^jH(c,h)}{j!},\quad
 d_0=\delta_1(c),\quad d_1=\delta_1'(c),\quad
 d_2=\tfrac12\delta_1''(c),\quad e_0=\delta_2(c).
\]
Then
\begin{equation}\label{eq:qfold-three}
 a=c+u_0Q^{1/2}+u_1Q+u_2Q^{3/2}+O(Q^2),
 \qquad u_0^2=\frac{Y-\delta_1(c)}{\alpha_2}.
\end{equation}
In particular,
$\delta_1'=-2\pi\sin(2\pi a)$,
$\delta_1''=-4\pi^2\cos(2\pi a)$, and
\[
 \delta_2=\cos(2\pi a)-1+\frac{\cos(4\pi a)-1}{2}.
\]
The two signs of $u_0$ give the two sheets. Their common $Q$ coefficient
$u_1$ need not be the physical critical-point shift: asymmetry of the
cubic forward term also shifts the midpoint of the two roots.

The moving-fold representation \eqref{eq:fold-exact}, applied to
$\mathcal F$, is the correct expression when the target approaches the
physical branch point $v_F(h)$. Expansion \eqref{eq:qfold-three} alone
is not uniform at that additional coalescence.

\subsection{A genuine critical Stokes comparison}
The preceding real example compares a physical completion with the
median-forward sum. To see the same action division in an actual
lateral Stokes comparison, let $c_-(h)$ be the critical point of $B_-$
near $a_*$, put $v_-=B_-(c_-,h)$, and set
$\alpha_-=(B_-)_{aa}(c_-,h)/2$. Since
$B_+=B_-+J(a,Q)$ and
$J=2\ii\sin(2\pi a)Q+O(Q^2)$, the upper lateral function has, at the
lower lateral critical target, the two local inverse values
\begin{equation}\label{eq:actual-stokes-fold}
 G_{+,\pm}(v_-,h)=c_-(h)\ \pm\
 \sqrt{-\frac{2\ii\sin(2\pi c_-(h))}{\alpha_-(h)}}\,
 Q^{1/2}+O(Q).
\end{equation}
The inverse of $B_-$ itself has a double value $c_-$ at that target.
The coefficient in \eqref{eq:actual-stokes-fold} is nonzero for small
$h$, since $a_*\in(1,3/2)$. On neighboring target patches the full
comparison is \eqref{eq:critical-jump}, not a discontinuous rule for
choosing the sign of a square root.

The physical $q$-gamma function is analytic across the positive $h$
axis. Its completion terms in \eqref{eq:gamma-plus} and
\eqref{eq:gamma-minus} cancel the jump of the perturbative lateral sums.
Equation \eqref{eq:actual-stokes-fold} compares those lateral sums; it
does not claim that the physical function itself has a discontinuity.

\section{A convergent reflection model with no Stokes jump}
\label{sec:reflection}

A particularly transparent critical example is obtained without any
divergent algebraic block. Consider
\begin{equation}\label{eq:reflection-def}
 R(a,h)=\log\bigl(\Gamma_{\e^{-h}}(a)
                         \Gamma_{\e^{-h}}(1-a)\bigr),\qquad 0<a<1.
\end{equation}
The identity $B_{2k+1}(1-a)=-B_{2k+1}(a)$ cancels the entire tail.
The ordinary gamma reflection formula \cite{dlmfgamma} gives the exact convergent core
\begin{equation}\label{eq:reflection-core}
 R_0(a,h)=\log\frac{\pi}{\sin(\pi a)}
       +\Log\frac{1-\e^{-h}}h+\frac{a(1-a)}2h.
\end{equation}
The completion and jump are
\begin{equation}\label{eq:reflection-completion}
 R(a,h)=R_0(a,h)+2D(a,Q),\qquad
 J(a,Q)+J(1-a,Q)=0.
\end{equation}
Thus $R_0$ is genuinely convergent near $h=0$; it has no lateral Borel
ambiguity at all. Nevertheless, the flat term $2D$ remains.

\begin{theorem}[Explicit reflection inverse transseries]
\label{thm:reflection}
The equation
\begin{equation}\label{eq:reflection-target}
 R(a,h)=R_0(1/2,h)
\end{equation}
has two inverse branches near $a=1/2$ of the form
\begin{equation}\label{eq:reflection-exact}
 a_\pm=\frac12\pm Q^{1/2}U(h,Q),
\end{equation}
where $U$ is jointly holomorphic near $(h,Q)=(0,0)$, and
\begin{equation}\label{eq:reflection-U0}
 U(h,0)=\sqrt{\frac8{\pi^2-h}}.
\end{equation}
The square root is positive at $h=0$. In particular, the complete
lifted transseries converges in $(h,Q)$, and only odd powers of
$Q^{1/2}$ occur in the displacement from $1/2$.
\end{theorem}
\begin{proof}
Both $R_0(1/2+z,h)-R_0(1/2,h)$ and $D(1/2+z,Q)$ are even holomorphic
functions of $z$, so they are holomorphic in $w=z^2$. Their equation is
\begin{equation}\label{eq:reflection-w}
 \frac{\pi^2-h}{2}w+O(w^2)-4Q+O(Qw,Q^2)=0.
\end{equation}
At $(w,h,Q)=(0,0,0)$ the $w$-derivative is $\pi^2/2\ne0$.
The analytic implicit theorem gives a unique holomorphic $w(h,Q)$
vanishing at $Q=0$. It factors as $w=QV(h,Q)$ with
$V(h,0)=8/(\pi^2-h)$, which is nonzero near $h=0$.
Choose $U=\sqrt V$. Then $z=\pm\sqrt Q\,U$ proves the theorem.
\end{proof}
This example isolates two phenomena that are often conflated. A Stokes
jump is absent, but a flat completion is present; that completion still
creates half-action inverse roots. It also shows that divergent
perturbative series are not necessary for exponentially amplified
critical inverse errors.

\subsection{Explicit higher sectors}
Write $w=b_1Q+b_2Q^2+b_3Q^3+\cdots$ and set
\[
 \alpha=\frac{\pi^2-h}{2},\qquad
 \beta_4=\frac{\pi^4}{12},\qquad\beta_6=\frac{\pi^6}{45}.
\]
Expansion of \eqref{eq:reflection-w} gives
\begin{align}
 b_1&=\frac4\alpha,\label{eq:rb1}\\
 b_2&=\frac{4-\beta_4b_1^2-4\pi^2b_1}{\alpha},\label{eq:rb2}\\
 b_3&=\frac{\frac{16}{3}-2\beta_4b_1b_2-\beta_6b_1^3
            -4\pi^2b_2+\frac43\pi^4b_1^2+4\pi^2b_1}{\alpha}.
            \label{eq:rb3}
\end{align}
Consequently,
\begin{equation}\label{eq:reflection-three}
 a_\pm=\frac12\pm\sqrt{b_1}Q^{1/2}
 \left[1+\frac{b_2}{2b_1}Q+
 \left(\frac{b_3}{2b_1}-\frac{b_2^2}{8b_1^2}\right)Q^2+O(Q^3)\right].
\end{equation}
Every coefficient is an analytic function of $h$ near zero. At $h=0$,
$b_1=8/\pi^2$ and $b_2/(2b_1)=-25/6$. These are coefficients of the
lift, not a substitution of $h=0$ into the original essential variable
$Q=\e^{-4\pi^2/h}$.

For $0<h<\pi^2$, $R_0$ is strictly convex on $(0,1)$ because its
second derivative is $\pi^2\csc^2(\pi a)-h>0$. The physical $R$ is
strictly convex there by \eqref{eq:Fsecond}. Both are symmetric about
$1/2$, while the physical minimum is lower by
\[
 2D(1/2,Q)=-4Q-4Q^2+O(Q^3).
\]
Hence the two real roots of \eqref{eq:reflection-target} are not an
artifact of complex branch notation. The omitted flat term changes the
actual inverse multiplicity at the chosen real target.

\section{Accuracy, action budgets, and compatibility of charts}
\label{sec:accuracy}

\subsection{Critical residuals are H\"older, not Lipschitz}
\begin{lemma}[Critical root-cluster certificate]\label{lem:critical-bound}
Suppose $F_0(c+z)-v=z^r A(z)$ on $|z|\le R_0$, with
$|A(z)|\ge a_0>0$. If $|F(c+z)-F_0(c+z)|\le\eps$ there, and
$R=(2\eps/a_0)^{1/r}<R_0$, then $F(x)=v$ has exactly $r$ roots,
counting multiplicity, in $|x-c|<R$.
\end{lemma}
\begin{proof}
On $|z|=R$, $|F_0(c+z)-v|\ge a_0R^r=2\eps$.
Rouch\'e's theorem compares $F-v$ with $z^rA(z)$, which has exactly
$r$ roots in that disk.
\end{proof}
This gives root existence and the scale $\eps^{1/r}$. It does not by
itself label the roots. After the scaled roots separate, the ordinary
disk certificate applies in the variable $u$ from
Theorem~\ref{thm:critical} and supplies branchwise errors.

An inverse tolerance on the scale $\e^{-B/h}$ at an order-$r$ critical
point generally demands a forward tolerance on the scale
$\e^{-rB/h}$. Every fixed algebraic remainder $O(h^M)$ is much larger
than the $Q$-sized target window in the $q$-gamma example. Substituting a
few more algebraic terms is therefore not a substitute for identifying
the flat completion. Even an optimally truncated approximation must
carry a sufficiently accurate remainder description to determine the
leading root coefficient in that window.

\subsection{More general monomials and a small curvature}
Theorem~\ref{thm:critical} is an analytic parameter theorem and can be
applied to any chosen small holomorphic monomial
\[
 t(h)=h^\beta(\Log h)^m\e^{-\mathcal A(h)}
\]
on a fixed cover. The inverse monomial is $t(h)^{1/r}$ when its leading
coefficient is nonzero and the uniform hypotheses hold. Thus algebraic,
logarithmic, and nonlinear exponential actions can all be transported
by the same local construction.

The nonvanishing curvature hypothesis matters. For example,
\[
 F(x,h)=\e^{-B/h}x^r+\e^{-A/h},\qquad A>B>0,
\]
at target zero has roots of size $\e^{-(A-B)/(rh)}$, not
$\e^{-A/(rh)}$. Dividing by the small leading coefficient changes the
effective action before ramification. An analogous normalization is
valid for more general cores only when the normalized higher derivatives
remain controlled on the required shrinking chart.

\subsection{Three distinct sources of multivaluedness}
There are three independent continuation operations in the examples.
First, changing the logarithm of $P$ changes $\ell$ by a multiple of
$2\pi\ii$ in \eqref{eq:h-inverse-phase}. Second, continuing around a
critical target value permutes the square-root or higher Puiseux
sheets. Third, changing a Borel summation direction across a singular
ray changes a sectorial block by its Stokes map. Each has a different
base variable and a different transition law.

A useful analytic atlas therefore records a source sector, a target
patch, all logarithm branches, the selected inverse sheet, and the
summation/completion convention. On overlaps, equations
\eqref{eq:target-transition}--\eqref{eq:source-transition} give the
regular transition identities; the moving fold chart supplies the
branched extension. This is a local compatibility framework, not a
claim of a single globally ordered transseries field on the entire
complex plane.

\subsection{What is and is not covered at other \texorpdfstring{$q$}{q}-critical points}
The explicit proofs concern $q\to1$ with $z=q^a\to1$. For a primitive
$k$th root $\zeta$, the exact residue-class splitting
\begin{equation}\label{eq:root-split}
 (\e^{-ah};\zeta\e^{-h})_\infty
 =\prod_{j=0}^{k-1}
   (\zeta^j\e^{-(a+j)h};\e^{-kh})_\infty
\end{equation}
separates a moving gamma block from the nontrivial cyclotomic blocks.
The factor $j=0$ is covered here after $h\mapsto kh$, $a\mapsto a/k$.
The other factors require additional cyclotomic completion data and may
interfere after differentiation. We do not transfer the action
$4\pi^2$ or the median criterion unchanged to that entire product.
The preceding repository q-cusp report treats a different modular
Euler-quotient subclass \cite{prior}.

\section{Reproducible verification and numerical diagnostics}
\label{sec:verification}

\subsection{Exact checks}
The accompanying \texttt{verify.py} uses SymPy for finite exact identities
and mpmath for high-precision diagnostics. It requires neither repository
files nor network access. The command
\begin{quote}\ttfamily
python verify.py --output results.json
\end{quote}
reproduces the recorded tests. The supplied run uses 110 decimal digits
and passed 144 exact assertions. These include the Bernoulli generating
kernel through degree $12$, the first gamma coefficients, finite-product
$q$-gamma expansions at integers $2$ through $8$ through degree $6$,
the second-order inverse response, the fold coefficients
\eqref{eq:u1}--\eqref{eq:u2}, the polynomial--logarithmic inverse,
the three reflection sectors, and finite M\"obius inversions.

Two exact rational comparisons also justify the constants used to
locate $a_*$ in \eqref{eq:astar}. Since
$\gamma<H_{30}-\log30$ and
\[
 \log30=2\sum_{j\ge0}\frac{(29/31)^{2j+1}}{2j+1},
\]
retaining $100$ positive terms gives an exact lower bound for $\log30$
that makes $H_{30}-\log30<3/5$. Also the first six terms of
$\exp(7/10)$ already exceed $2$, proving $\log2<7/10$.
These tests are finite rational calculations, not floating-point
estimates of the endpoint signs.

\subsection{A numerically independent completion check}
To compare the physical product with the Borel calculation, the program
evaluates $p(a,h)$ directly using the convergent $q$-product. It evaluates
$\Med\widehat\Phi_a$ separately by the kernel series
\eqref{eq:median-kernel}. The latter is accelerated by subtracting its
first ten algebraic kernel terms and summing those subtractions exactly
through the known Bernoulli coefficients:
\begin{align}
 \Med\widehat\Phi_a(h)
 ={}&\sum_{k=1}^{K}c_{2k}(a)h^{2k}\notag\\
 &+\frac2\pi\sum_{n\ge1}\sigma_n(a)
 \left[K(An/h)-\sum_{k=1}^{K}\frac{(2k-1)!}{(An/h)^{2k}}\right],
 \label{eq:acceleration}\\
 c_{2k}(a)&=\frac{B_{2k}B_{2k+1}(a)}{2k(2k+1)!},\qquad
 \sigma_n(a)=\sum_{d\mid n}\frac{\sin(2\pi da)}d.\notag
\end{align}
This identity follows by absolute convergence after subtraction and the
Bernoulli Fourier identities. The finite numerical cutoff is not
outward-rounded or interval-certified. Comparing action cutoffs $60$
and $120$ is a diagnostic, not a certified tail enclosure.

\begin{center}
\begin{tabular}{@{}cccc@{}}
\toprule
$a$ & $h$ & Completion residual, cutoff $120$ & Change from cutoff $60$\\
\midrule
$0.3$ & $2$ & $8.03\times10^{-55}$ & $1.55\times10^{-48}$\\
$0.5$ & $2$ & $8.32\times10^{-112}$ & $5.07\times10^{-159}$\\
$0.7$ & $3$ & $6.01\times10^{-51}$ & $1.16\times10^{-44}$\\
\bottomrule
\end{tabular}
\end{center}
The tiny middle row is a degenerate exact-tail case:
$\widehat\Phi_{1/2}=0$. It is not evidence of superior numerical
conditioning for general $a$. The residuals verify the chosen signs and
normalizations to the displayed numerical accuracy. The proof of the
infinite identity is the Mellin argument in Section~\ref{sec:mellin}.

\subsection{Critical-root scaling}
For the natural gamma minimum, the program constructs $H=F-D$ from
independent physical and dual-product formulas, finds its critical point
$c$, and solves the physical equation $F(a,h)=H(c,h)$. The root solve
is performed in the scaled variable
\[
 u=\frac{a-c}{\sqrt{(1-\cos(2\pi c))/\alpha}\,\sqrt Q},
\]
not in the ill-conditioned unscaled variable. The predicted limiting
values are $u=\pm1$.

\begin{center}
\begin{tabular}{@{}cccc@{}}
\toprule
$h$ & $Q$ & Predicted positive displacement & $|u_+-1|$\\
\midrule
$2$ & $2.67529\times10^{-9}$ & $1.79327\times10^{-4}$ & $2.292\times10^{-4}$\\
$1$ & $7.15717\times10^{-18}$ & $6.98150\times10^{-9}$ & $6.011\times10^{-9}$\\
$0.5$ & $5.12250\times10^{-35}$ & $1.63654\times10^{-17}$ & $1.122\times10^{-17}$\\
\bottomrule
\end{tabular}
\end{center}
The two physical root residuals in these tests are below
$2\times10^{-111}$ at the working precision. This is a floating-point
residual, not an interval-certified bound on the exact roots.
The displacement is much larger than $Q$, and the normalized error
decreases on the expected $\sqrt Q$ scale.

The reflection model is checked against the \emph{actual} product
$\Gamma_q(a)\Gamma_q(1-a)$ rather than only its completion formula.
At $h=1$, the predicted leading displacement is
$2.54075862518576\times10^{-9}$. The observed coefficient
$(u_+-1)/Q$ is $-4.7764454296641167\ldots$, versus
$b_2/(2b_1)=-4.7764454296641171\ldots$. At $h=2$, the corresponding
values differ by about $2.23\times10^{-7}$, consistent with the next
$Q$ correction. Full values are retained in \texttt{results.json}.

\subsection{What was not verified computationally}
No computation proves the Mellin contour deformation, a resurgent
closure theorem, the Banach-analytic parameter result, or a universal
complex continuation theorem. Those are addressed by the conventional
proofs and the stated classical dependencies. No Lean proof, interval
root certificate, or independently refereed verification is claimed.
The tests are intended to make sign, coefficient, and branch mistakes
easier to detect and reproduce.

\section{Further research questions and topics}
\label{sec:further}

The following questions are not asserted to be solved by this article.
They are organized around specific obstructions rather than a claim of
unrestricted inversion of every possible transseries.

\subsection{Inverse Borel singularities and local monodromy}
For the regular inverse of $a\mapsto\log\Gamma_q(a)$, determine the
precise Borel singularities of each inverse coefficient block, not just
the allowed additive action lattice. The forward kernel is meromorphic
with explicitly known residues. Nonlinear inversion introduces
convolutions and can change poles into logarithmic or more complicated
singularities. A useful goal is an explicit bridge equation for the
inverse in a selected target chart, together with its continuation when
the target approaches a fold.

\subsection{A full nonlinear median at a critical value}
The arithmetic median used here is a linear logarithmic construction.
Build an algebra-compatible transseries completion for the full inverse
family, with its transseries parameters transformed by Stokes maps.
Determine how its median convention interacts with the ramified
coordinate $\sqrt{y-v(h)}$. Proposition~\ref{prop:median} shows that
averaging already-inverted lateral branches is not an acceptable shortcut.

\subsection{Confluence with gamma poles and zeros}
The present parameter domains remain away from the nonpositive integers.
When $a$ approaches a pole of $\Gamma(a)$ simultaneously with $h\to0$,
the gamma core itself becomes singular. Construct a uniform inverse
chart for $a=-m+\lambda h^\nu$ and identify the correct gamma,
logarithmic, and exponential scales. The local continuation of the
Borel kernel is not enough: pole removal and inverse-sheet selection
must be controlled together.

\subsection{Cyclotomic completions at general roots of unity}
Use \eqref{eq:root-split} to extend the Mellin argument to the complete
residue-class family. The Fourier functional equation is then replaced
by a finite phase-coupled system. Determine its exact even completion,
its lateral jumps, and the analogue of Theorem~\ref{thm:weights}.
One should expect the inverse critical geometry to depend on cancellations
between cyclotomic blocks, not merely on their individual action sizes.

\subsection{Higher critical points in natural gamma quotients}
Study finite products
\[
 \prod_{j=1}^{m}\Gamma_q(\lambda_j a+\mu_j)^{r_j}
\]
on domains where all limiting gamma arguments are regular. Their
completion is an explicit linear combination of $D$, and their jump is
an explicit combination of $J$. Find natural parameter families in
which the first $r-1$ derivatives of the limiting logarithm vanish,
while the $r$th does not. Theorem~\ref{thm:critical} then predicts an
$A/r$ action, unless a completion jet cancels. Classify when those
cancellations are forced by reflection or multiplication identities.

\subsection{Joint action and algebraic optimal truncation}
An inverse approximation has two distinct truncation indices: the order
inside an algebraic block and the number of exponential blocks retained.
Develop certified joint choices for a prescribed inverse tolerance,
using the exact kernel and the critical root-cluster certificate.
A particularly useful target is a rigorous error bound that remains
uniform while $y-v_F(h)$ passes through the $Q$ window and then through
a smaller window where the two inverse roots are close to merging.

\subsection{Countable actions beyond local finiteness}
Proposition~\ref{prop:countable} gives a Banach-analytic inverse for
normally summable perturbations. Determine which of these inverses also
possess a canonical formal transseries representation when action
values accumulate. A norm can make an analytic sum meaningful even
when a usual coefficient-by-coefficient action filtration fails.
The missing ingredient is a support-sensitive uniqueness theorem,
not just another contraction estimate.

\subsection{Global inverse atlases and growing cusp denominator}
For $q$-families continued between rational boundary charts, construct
source and target atlases that track logarithmic sheets, Stokes maps,
and critical-value monodromy simultaneously. Then investigate joint
limits in which a root-of-unity denominator grows with $1/h$.
The local fixed-denominator identities do not automatically provide
uniform separation of inverse sheets or a globally summable action set.

\subsection{Basic hypergeometric sums and saddle interference}
Extend the completed-product method to basic hypergeometric functions,
Nahm sums, and products occurring in quantum-dilogarithm integrals.
Their forward asymptotics may contain several saddles rather than one
gamma core. Derivative cancellation between saddles can create inverse
critical points on a different scale from the forward Stokes transition.
A successful theorem should identify both the saddle completion and the
inverse branch before claiming an error bound.

\subsection{Formalization and certification}
A practical formalization can proceed in layers. The finite Fourier and
M\"obius criterion, the inverse coefficient identities, and the
rational sign bounds are algebraic starting points. The next layer is
the filtered contraction and the analytic fold factorization. The
Mellin--Borel theorem then needs library support for principal-value
kernels, contour shifts, and uniform infinite sums. Separately, a
complex-interval implementation could certify particular inverse roots
using Lemmas~\ref{lem:certificate} and \ref{lem:critical-bound} without
first formalizing the full resurgence theory.

\section{Conclusion}
The missing information in a complex asymptotic inverse is often not
another algebraic coefficient. It is the analytic completion and the
location of the moving inverse branch point. For the moving
$q$-Pochhammer family, the Mellin calculation separates the remainder
into a sine part, which is the symmetric Borel sum, and a cosine part,
which is an independent convergent dual product. This proves the two
specified summability statements and makes their scope explicit.

Once the completed function is known, regular analytic inversion gives
full exponential sectors, while critical scaling gives fractional
actions. The $q$-gamma minimum exhibits both a genuine half-action
Stokes comparison and a real physical inverse splitting. The reflection
product gives an even cleaner limiting lesson: a convergent core and
zero Stokes jump can coexist with a nonzero flat term that changes the
inverse problem at leading exponential order.

The resulting calculus is local, branch-aware, and exact in its flat
parameters. It does not eliminate the need to choose a summation class,
a sector, or an inverse sheet; it makes those choices part of the
mathematical statement.

\appendix
\section{Proof dependencies and research-status audit}
\label{app:audit}

\begingroup\small
\begin{longtable}{@{}P{.24\textwidth}P{.34\textwidth}P{.36\textwidth}@{}}
\toprule
Result & Main ingredients & Boundary of the claim\\
\midrule
\endhead
Borel kernel and residues & Bernoulli generating functions; normally
convergent meromorphic sums & The rational forward singularity data are
already known; complex-parameter continuation uses the non-Fourier
kernel.\\[5pt]
Exact completion & Mellin inversion; Hurwitz and Riemann functional
equations; principal-value kernel & First proved for $h>0$, $0<a<1$;
then analytic continuation in the specified domains. No residue-only
argument is substituted for reconstruction.\\[5pt]
Conjecture 5 comparison & Explicit rotation $h=-2\pi\ii Ny$ and removal
of the dual index-zero factor & $1\le k<N$, $\Im y>0$, with the specified
lateral logarithms. No value is assigned to $\log(1-1)$.\\[5pt]
Weight criterion and Conjecture 1 & Completion identity; divisor
M\"obius inversion; finite Fourier inversion & Arbitrary finite complex
weights, including characters. This is a linear median statement.\\[5pt]
Regular inversion & Classical observable Lagrange inversion; analytic
implicit theorem; disk contraction & A selected regular chart. Summation
compatibility requires a suitable summable algebra.\\[5pt]
Critical action division & Ramified analytic blow-up and simple scaled
roots & Compact target patches away from the scaled discriminant;
nondegenerate normalized leading coefficient.\\[5pt]
Moving-fold chart & Parametric analytic Morse factorization & Valid at
the moving branch point on a square-root cover, not as a single-valued
holomorphic inverse in the target.\\[5pt]
Natural gamma fold & Exact completion; positive $q$-trigamma; uniform
parameter asymptotics & Small $h$ near the classical positive minimum.
All exponential estimates keep the moving coefficient blocks.\\[5pt]
Reflection inverse & Gamma reflection; Bernoulli parity; analytic IFT
in $z^2,h,Q$ & A genuinely convergent local lift. The physical variable
relation is imposed only after constructing the lift.\\[5pt]
Numerical results & Exact symbolic checks and 110-digit evaluation &
Not interval certificates, independent referee reports, or
proof-assistant verification.\\
\bottomrule
\end{longtable}
\endgroup

\subsection*{Scope of the source audit}
The repository comparison uses its transseries documentation and the
identified prior report's fixed-parameter, jump-transport, and further
research sections. It is not a statement-by-statement audit of the full
canonical volume. The cited public literature was checked on 4 October
2026. The inspected Fantini--Rella version explicitly labels the two
summability statements as conjectures; the present proof targets their
formulas rather than a broader slogan about modular resurgence.

The derived Mellin argument and applications have undergone an internal
algebraic and analytic consistency review in preparing this package.
The principal-value interchange, lateral orientation, dual product
constant, complex-parameter continuation, and critical target exclusions
are spelled out to facilitate external review. Such an internal review
cannot certify first publication priority or replace an independent
expert assessment.

\section{A compact formula and sign reference}
\label{app:signs}
The conventions most likely to cause a sign error are collected here.
Let $\Phi_\pm=\mathcal S_\pm\widehat\Phi_a$,
$L_\pm=\log(\e^{\pm2\pi\ii a}Q;Q)_\infty$, and
$L_0=\log(Q;Q)_\infty$. Then
\begin{align*}
 J&=\Phi_+-\Phi_-=L_--L_+,
       &J&=2\ii\sin(2\pi a)Q+O(Q^2),\\
 \mathcal S_+\widehat p&=E-\Phi_+=p-L_-,
       &\mathcal S_-\widehat p&=E-\Phi_-=p-L_+,\\
 p-\Med\widehat p&=(L_++L_-)/2,
       &p-\Med\widehat p&=-\cos(2\pi a)Q+O(Q^2),\\
 F-B_+&=L_0-L_-,
       &F-B_-&=L_0-L_+,\\
 F-H&=D=L_0-(L_++L_-)/2,
       &D&=(\cos(2\pi a)-1)Q+O(Q^2).
\end{align*}
The contour difference ``upper minus lower'' is clockwise and equals
$-2\pi\ii$ times the residue. This explains the positive sign in the
first line despite the negative residue in \eqref{eq:residues}.

For a regular inverse of $f+td$, the first response is $-td/f'$.
For the $h$-inverse of the factorial, the first physical correction is
therefore positive when $\cos(2\pi a)>0$, as displayed in
\eqref{eq:24-h}. For the gamma parameter inverse the first correction
is $(1-\cos(2\pi a_H))Q/H_a$, as in \eqref{eq:a-flat}.
At a fold these linear expressions are replaced by the ramified balance.

\begin{thebibliography}{99}
\raggedright

\bibitem{repo}
V. Reshetnikov, \emph{ProveIt}, repository documentation, snapshot
\texttt{b484725c893a7a9bde856151988996aef6965cb0}, inspected 4 October 2026.
Relevant path:
\path{Analysis/Transseries/docs/series-and-transseries/README.md}.
\url{https://github.com/VladimirReshetnikov/ProveIt}.

\bibitem{prior}
ProveIt research package, \emph{Complex Transseries Reversion at
$q$-Cusps: Convergent normal forms, infinitely flat limits,
logarithmic--Puiseux inverses, and recovery from cusp data}, 4 October
2026. Prior user-library article, TeX file
\texttt{article(20261004-191746).tex}; its stated repository snapshot is
\texttt{8e9cd6f00e0ac79c4425ff6abdfa64b497cf3f41}.
Consulted sections on shifted factorials, sectorial transport, and
further research. Unrefereed research report.

\bibitem{fr}
V. Fantini and C. Rella, \emph{Modular resurgence, $q$-Pochhammer
symbols, and quantum operators from mirror curves}, Letters in
Mathematical Physics \textbf{116} (2026), article 39.
\url{https://doi.org/10.1007/s11005-026-02063-x}.
Version specifically compared: arXiv:2506.08265v2, 1 April 2026,
Conjectures 1 and 5, formulas (60a)--(60b), and Theorem 4.4.
\url{https://arxiv.org/html/2506.08265v2}.

\bibitem{ar}
A. Arabi Ardehali and H. Rosengren, \emph{A new product formula for
$(z;q)_\infty$, with applications to asymptotics}, arXiv:2602.11329,
2026. In particular, the moving-parameter expansion in Corollary 3.
\url{https://arxiv.org/abs/2602.11329}.

\bibitem{mcintosh}
R. J. McIntosh, \emph{Some asymptotic formulae for $q$-shifted factorials},
The Ramanujan Journal \textbf{3} (1999), 205--214.
\url{https://doi.org/10.1023/A:1006949508631}.

\bibitem{sauzin}
D. Sauzin, \emph{Nonlinear analysis with resurgent functions},
Annales scientifiques de l'\'Ecole Normale Sup\'erieure (4)
\textbf{48} (2015), no.~3, 667--702.
\url{https://doi.org/10.24033/asens.2255}.
Preprint: \url{https://arxiv.org/abs/1212.4477}.

\bibitem{sauzinintro}
D. Sauzin, \emph{Introduction to 1-summability and resurgence},
arXiv:1405.0356, 2014.
\url{https://arxiv.org/abs/1405.0356}.

\bibitem{dlmfzeta}
NIST Digital Library of Mathematical Functions,
\emph{\S25.11: Hurwitz Zeta Function}, especially (25.11.9),
(25.11.13), (25.11.14), and (25.11.18).
\url{https://dlmf.nist.gov/25.11}. Accessed 4 October 2026.

\bibitem{dlmfriemann}
NIST Digital Library of Mathematical Functions,
\emph{\S25.4: Reflection Formulas for the Riemann Zeta Function}.
\url{https://dlmf.nist.gov/25.4}. Accessed 4 October 2026.

\bibitem{dlmfqgamma}
NIST Digital Library of Mathematical Functions,
\emph{\S5.18: $q$-Gamma and $q$-Beta Functions}.
\url{https://dlmf.nist.gov/5.18}. Accessed 4 October 2026.

\bibitem{dlmfgamma}
NIST Digital Library of Mathematical Functions,
\emph{\S5.4: Special Values and Extrema}, and
\emph{\S5.5: Functional Relations}, including the gamma reflection
formula and the special digamma values used here.
\url{https://dlmf.nist.gov/5.4}; \url{https://dlmf.nist.gov/5.5}.
Accessed 4 October 2026.

\end{thebibliography}
\end{document}
